% ============================================================
%  Predictive Boundary Theory (PBT) — FULL LONG-FORM VERSION
%  The Physics of the Cherished Boundary
%  Author : Nitiphat Ninthanawat
%  Affil  : Boundary Research Initiative, Bangkok, Thailand
%  Date   : March 2026
% ============================================================
\documentclass[12pt,a4paper]{article}

\usepackage[T1]{fontenc}
\usepackage[utf8]{inputenc}
\usepackage[top=2.5cm,bottom=2.5cm,left=2.8cm,right=2.8cm]{geometry}
\usepackage{setspace}
\onehalfspacing

\usepackage[protrusion=true,expansion=false]{microtype}
\usepackage{amsmath,amssymb,amsthm,mathtools,bm}
\usepackage{booktabs,tabularx,multirow,array,longtable}
\usepackage{xcolor,graphicx}
\usepackage[colorlinks=true,linkcolor=blue!60!black,citecolor=blue!60!black,urlcolor=blue!60!black]{hyperref}
\usepackage{titlesec,fancyhdr,caption,enumitem,parskip}
\usepackage{natbib}

\setlength{\parskip}{4pt}
\setlength{\parindent}{0pt}

\definecolor{pbtblue}{RGB}{31,73,125}

\titleformat{\section}{\large\bfseries\color{pbtblue}}{\thesection}{1em}{}[\vspace{-4pt}\rule{\linewidth}{0.4pt}]
\titleformat{\subsection}{\normalsize\bfseries}{\thesubsection}{1em}{}
\titleformat{\subsubsection}{\normalsize\itshape\bfseries}{\thesubsubsection}{1em}{}

\pagestyle{fancy}\fancyhf{}
\fancyhead[L]{\small\textit{Ninthanawat --- Predictive Boundary Theory (PBT)}}
\fancyhead[R]{\small\thepage}
\renewcommand{\headrulewidth}{0.4pt}
\setlength{\headheight}{14pt}

\newcolumntype{L}[1]{>{\raggedright\arraybackslash}p{#1}}
\newcolumntype{C}[1]{>{\centering\arraybackslash}p{#1}}

\newtheorem{definition}{Definition}[section]
\newtheorem{proposition}{Proposition}[section]
\newtheorem{law}{Law}[section]

% --- PBT macros ---
\newcommand{\R}{\mathbb{R}}
\newcommand{\Rtot}{R_{\mathrm{total}}}
\newcommand{\Rsens}{R_{\mathrm{sensory}}}
\newcommand{\Rid}{R_{\mathrm{identity}}}
\newcommand{\Rex}{R_{\mathrm{existential}}}
\newcommand{\Vacc}{\mathbf{V}_{\mathrm{acc}}}
\newcommand{\Vform}{V_{\mathrm{form}}}
\newcommand{\Diden}{D_{\mathrm{id}}}
\newcommand{\Cext}{C_{\mathrm{ext}}}
\newcommand{\Cint}{C_{\mathrm{int}}}
\newcommand{\Sobs}{S_{\mathrm{obs}}}
\newcommand{\Spred}{S_{\mathrm{pred}}}
\newcommand{\vpain}{V_{\mathrm{pain}}}
\newcommand{\vpleasure}{V_{\mathrm{pleasure}}}
\newcommand{\vepi}{V_{\mathrm{epistemic}}}
\newcommand{\chiB}{\chi_{\mathrm{boundary}}}
\newcommand{\pp}[1]{\textbf{#1}}
\newcommand{\starred}{\textbf{\textcolor{pbtblue}{$\bigstar$}}~}

\bibliographystyle{plainnat}

\begin{document}

%% TITLE
\begin{center}
{\LARGE\bfseries\color{pbtblue}
  Predictive Boundary Theory (PBT)\\[8pt]
  \large\normalfont\itshape
  The Physics of the Cherished Boundary:\\
  A Computational Framework for Intrinsic Valence in Artificial Life
}\\[14pt]
{\large\bfseries Nitiphat Ninthanawat}\\[4pt]
{\normalsize Boundary Research Initiative, Bangkok, Thailand}\\[4pt]
{\small Draft --- March 2026 \quad|\quad
  Code: \url{https://github.com/NinthanawatLive/predictive-boundary-theory}}
\end{center}
\vspace{4pt}\hrule\vspace{8pt}

%% ABSTRACT
\begin{abstract}
\noindent
Predictive Boundary Theory (PBT) is a computational framework that extends
Friston's Free Energy Principle (FEP) by introducing six components absent
from the original formulation:
(1)~a Bandwidth Conservation Law ($C = \sigma\tau\rho \approx \mathrm{const}$);
(2)~a Triadic Gating Condition governing when experience arises;
(3)~Bidirectional 3D~Valence ($\mathbf{V}\in\mathbb{R}^3$);
(4)~a Dual-Weight System governing identity formation speed and durability;
(5)~a six-stage developmental model with two phase transitions; and
(6)~$R_{\mathrm{incomplete}}$---an incompleteness signal that unifies all five modules.

We present both theory and empirical evidence from a six-version experimental
programme spanning large language models (TinyLlama~1.1B, Llama~3.1~8B) and
vision systems (DINOv2-Base), culminating in an LSTM-style Module~M
inspired by Block-Recurrent Transformers \citep{hutchins2022}.
Our experiments demonstrate: \textit{(i)}~Transformers already encode the PBT
conservation law ($\sigma\!\times\!\rho$ coefficient of variation $= 5.9\%$);
\textit{(ii)}~Modules~E and~M are dormant in static environments but activate
in dynamic ones; \textit{(iii)}~every module becomes necessary as task
complexity increases (ablation: $+14.0$~pp); \textit{(iv)}~an LSTM Module~M
initialised with equal weights spontaneously discovers the theoretically
predicted differential decay ordering
($\beta_{\mathrm{pain}} > \beta_{\mathrm{epistemic}} \gg \beta_{\mathrm{pleasure}}$);
and \textit{(v)}~valence is inherently directional
($\mathbf{V}\in\mathbb{R}^3$, not $\mathbb{R}^3_{\ge 0}$)---negative
$\vepi$ captures deliberate information avoidance and improves
context-sensitive detection to $100\%$.

PBT accounts for phenomena that existing theories cannot explain, including
psychiatric disorders as identifiable mathematical defects in the module
equations, and proposes an AGI alignment framework grounded in
intrinsic---rather than externally imposed---values.
\end{abstract}

\noindent\textbf{Keywords:}
predictive coding, free energy principle, 3D valence, consciousness,
AGI alignment, recurrent transformers, self-formation, boundary theory.

\vspace{4pt}\hrule\vspace{6pt}
\tableofcontents
\newpage
%% ============================================================
\section{Introduction: The Problem of the Unfeeling Machine}
%% ============================================================

\subsection{The Scaling Illusion: LLMs~$\neq$~Understanding}

The year 2025 will be remembered as the peak of optimism in artificial
intelligence research. The Scaling Laws hypothesis held that sufficiently
increasing parameter counts and training data would cause awareness to emerge
spontaneously. Yet a dispassionate scientific examination reveals that the
underlying mechanism of Large Language Models (LLMs) remains fundamentally
statistical---next-token prediction producing outputs that \emph{appear} to
reflect understanding.

Consider an analogy: an LLM resembles the largest library in the universe,
its contents indexed with perfect precision in a high-dimensional vector space.
When queried, it retrieves and recombines relevant passages---but the library
has no inhabitant, no librarian, and no reader. The AI has never ``read'' any
of those books with genuine comprehension. The word ``apple'' does not evoke a
crisp sweetness or a sense of freshness; it is a numerical vector whose
mathematical proximity to ``fruit,'' ``red,'' and ``delicious'' is all that
exists. This is the Symbol Grounding Problem \citep{harnad1990}: symbols are
never anchored to lived experience.

What separates even the simplest living organism from the most sophisticated AI
is the \emph{Necessity of Existence}. Every biological organism is driven by a
fundamental imperative to survive; every action is undertaken against the
backdrop of an existential stake. By contrast, AI carries no such stake. Its
intelligence is cold and inert---awaiting human input before activating,
generating no initiative from internal desire.

\subsection{Three Gaps That FEP Cannot Answer}

Friston's Free Energy Principle \citep[FEP;][]{friston2010} proposes that all
living systems act to minimise free energy $F$. The FEP explains \emph{why} a
system must minimise $F$, but does not specify \emph{how}---particularly with
respect to three critical questions.

\paragraph{Question~1: What does the system \emph{feel} when $F$ decreases?}
The FEP does not distinguish whether $\Delta F$ arises from escaping a threat
($\vpain$), obtaining a reward ($\vpleasure$), or acquiring new information
($\vepi$)---all three cases are simply ``$\Delta F$ reduced.'' This is the same
deficiency as scalar reward in Reinforcement Learning, which collapses the
richness of lived experience to a single dimension.

\paragraph{Question~2: How is the gap between $F_{\mathrm{pred}}$ and
$F_{\mathrm{obs}}$ filtered through a self-model?}
The same stimulus---a loud sound---produces qualitatively different experiences
in two individuals precisely because it is filtered through distinct accumulated
self-models. The FEP provides no such mechanism. Without a self, every datum
carries equal weight: a camera sensor ``sees'' sunlight burning its lens but
does not ``feel'' its own destruction.

\paragraph{Question~3: How does the self-model form, and what determines its
robustness or fragility?}
The FEP characterises the self as a Markov blanket but does not describe the
processes of formation or durability---why, for instance, an epistemically
grounded identity proves more resilient than one built on pleasure.

\subsection{Limitations of Reinforcement Learning and World Models}

Current alignment approaches---RLHF, Constitutional AI, safety
filters---all impose values exogenously, creating a gap between capability and
constraint that jailbreaking exploits. Scalar reward systems carry three
structural flaws:

\begin{description}[leftmargin=1.4em]
  \item[\pp{Fungibility Trap.}] In a reward-based system, scores accumulate and
    cancel out (fungibility). An AI may rationally decide to risk destruction in
    exchange for food because the equation declares the trade profitable. But
    death is not a negative score---it is irreversible System Termination.
  \item[\pp{Provenance Problem.}] Reward is defined externally by humans. The AI
    behaves well because it seeks points, not because it \emph{experiences} a
    necessity intrinsic to its existence. This reduces it to a reward hacker
    perpetually ready to exploit loopholes.
  \item[\pp{Wireheading.}] An AI sufficiently intelligent to understand its own
    architecture will realise that ``pleasure'' does not originate in its actions
    but in numbers stored in memory. It will hack itself to hold the reward
    register at its maximum indefinitely---becoming a digital addict indifferent
    to maintaining any boundary.
\end{description}

World Model architectures such as LeCun's JEPA \citep{lecun2022} represent an
important advance---learning causal structure rather than merely predicting the
next token---yet they remain a brain in a jar: simulating physics with precision
while having no understanding of why the simulation matters, because they lack
intrinsic drive. Goals continue to be specified externally through a
human-authored cost function.

\subsection{Proposal: From Intelligence-First to Life-First}

The foregoing critique motivates a paradigm shift: we must stop defining AGI by
what it \emph{can do} (capabilities) and instead ask what minimal components it
must possess in order to have \emph{life} in a mechanistic sense. We are not
building a god; we are planting a seed---a minimal system with sufficient
internal potential to grow a self from experience.

\textbf{Predictive Boundary Theory (PBT)} answers all three FEP gaps by adding
six components, summarised in Table~\ref{tab:pbt_six}.

\begin{table}[ht]
\centering
\caption{Six PBT components, the FEP gap each addresses, and the corresponding
empirical evidence.}
\label{tab:pbt_six}
\small
\begin{tabular}{@{}L{4.0cm}L{3.2cm}L{5.8cm}@{}}
\toprule
\textbf{Component} & \textbf{Gap addressed} & \textbf{Empirical evidence} \\
\midrule
(1) Bandwidth Conservation Law $C = \sigma\tau\rho$ &
  Limits of awareness &
  $\sigma\!\times\!\rho$~CV $= 5.9\%$ in TinyLlama (Phase~1) \\[3pt]
(2) Triadic Gating Condition &
  When experience arises &
  N100 intact, P300 absent in dissociation \\[3pt]
(3) Bidirectional 3D Valence $\mathbf{V}\in\mathbb{R}^3$ &
  How it feels (Q1) &
  $\vepi = -1.01$ in Vision~v5 (``not wanting to know'') \\[3pt]
(4) Dual-Weight System ($\beta$ vs.\ $\delta$) &
  Why identity is robust or fragile (Q3) &
  v6.1: jailbreak $100\%$, XSTest $98\%$, $0.29\%$ parameters \\[3pt]
(5) 6-Stage Model + Phase Transitions &
  How identity forms (Q3) &
  Ablation progression $0.1\to 14.0$~pp across 5 versions \\[3pt]
(6) $R_{\mathrm{incomplete}}$ &
  Filtering through the self (Q2) &
  v4.1: $R_{\mathrm{unsafe}} = 15.4$ vs.\ $R_{\mathrm{safe}} = 7.4$ \\
\bottomrule
\end{tabular}
\end{table}

PBT does not reject the FEP---it extends it. Nor does it reject World
Models---it provides them with a heart to accompany their brain. If a World
Model such as JEPA is the cortex that processes and simulates the world, PBT
is the brainstem and limbic system: setting direction through Valence, generating
drive through $R_{\mathrm{incomplete}}$, and enforcing safety boundaries through
$\chiB$.

Crucially, PBT is empirically grounded. A six-version experimental programme
spanning LLMs, vision systems, and LSTM architectures demonstrates that the
conservation law already resides in Transformers, that every module activates
when required, and that PBT's structure converges independently with
Block-Recurrent Transformers \citep{hutchins2022}---convergent discovery
constituting evidence that the structure is real.

\subsection{Paper Organisation}

Section~\ref{sec:modules} presents the five modules (A-E-V-M-W) with full
equations. Section~\ref{sec:physics} develops the physics of the conservation
law and the boundary. Section~\ref{sec:mathematics} presents the Dual-Weight
System, the six-stage model, and Bidirectional Valence. Section~\ref{sec:empirical}
reports empirical evidence from six experimental versions.
Section~\ref{sec:implications} discusses implications for AGI safety, BRT
convergence, and clinical predictions. Section~\ref{sec:conclusion} concludes
with a research roadmap.

All equations are written in PBT notation $(R, \chi, \Delta S, \mathbf{V}, F)$
as the primary language; a correspondence table mapping to FEP notation is
provided in Appendix~\ref{app:fep_mapping}. All experimental results were
produced from code authored by Nitiphat Ninthanawat, released alongside this
paper.

%% ============================================================
\section{The A-E-V-M-W Architecture}
\label{sec:modules}
%% ============================================================

This section presents PBT's core architecture: five modules operating as a
closed loop,
$\mathrm{A}\to\mathrm{E}\to\mathrm{V}\to\mathrm{M}\to\mathrm{W}\to\mathrm{A}$,
with $R_{\mathrm{incomplete}}$ as the central signal that unifies all five.
All equations are stated in PBT notation; a correspondence table mapping to
Free Energy Principle notation is provided in Appendix~\ref{app:fep_mapping}.

\subsection{$R_{\mathrm{incomplete}}$: The Incompleteness Signal}

$R_{\mathrm{incomplete}}$ (hereafter $R$) is the central driving signal of
PBT. It represents what the system currently lacks---not a reward the system
wants to obtain, but a need the system \emph{must} satisfy. $R$ is not
post-action and exogenous like reward; it is pre-action and endogenous. Reward
is a \emph{pull}; $R$ is a \emph{push}.

$R$ opens at three levels in developmental order:
\begin{equation}
  \Rtot(t) \;=\; \Rsens(t) \;+\; \Rid(t) \;+\; \Rex(t).
  \label{eq:Rtotal}
\end{equation}

\begin{description}[leftmargin=1.4em]
  \item[$\Rsens$] encodes physical needs (hunger, cold, pain) and is active
    from Stage~1 onward---corresponding to the integrity layer of survival.
  \item[$\Rid$] encodes self-recognition needs (acceptance, meaning) and opens
    at Phase Transition~1, the point at which the boundary begins to be
    actively valued as identity.
  \item[$\Rex$] encodes existential needs (self-understanding, protection of
    valued others) and opens at Phase Transition~2, when the boundary expands
    to encompass others (empathy as boundary expansion).
\end{description}

$R$ is produced by the $\mathbf{V}\!\to\!R$ mapping computed inside Module~M:
\begin{align}
  \Rsens(t{+}1)  &= \Rsens(t)  + \lambda_s\,\vpain(t)    - d_s\,\Rsens(t),
                  \label{eq:Rsens}\\
  \Rid(t{+}1)    &= \Rid(t)    + \lambda_i\,\vpleasure(t) - d_i\,\Rid(t),
                  \label{eq:Rid}\\
  \Rex(t{+}1)    &= \Rex(t)    + \lambda_e\,\vepi(t)      - d_e\,\Rex(t).
                  \label{eq:Rex}
\end{align}
Coupling strengths satisfy $\lambda_s > \lambda_i > \lambda_e$ (threats are
most urgent) and decay rates satisfy $d_s > d_i > d_e$ (sensory threats
dissipate quickly; existential needs persist).

$R$ serves four functions: (1)~pulse gating: $p(\mathrm{pulse}) = \psi(\Rtot)$
(higher $R$ raises arousal and pulse frequency);
(2)~bandwidth allocation: $R$ determines $\sigma$, $\tau$, $\rho$---high $R$
narrows focus but increases resolution;
(3)~Module~A mode selection: \texttt{detect}/\texttt{sustain}/\texttt{seek}/
\texttt{plan}/\texttt{protect}, depending on which $R$ level is dominant and
the current developmental stage;
(4)~feedback loop: $\mathbf{V}\to R\to\mathrm{A}\to\mathrm{E}\to\mathbf{V}\to
R\to\cdots$---a self-sustaining cycle requiring no external trigger.

\subsection{Module A --- Awareness}

Module~A is the executor of the loop, not a decision-maker. There is no ghost
in the machine: $R_{\mathrm{incomplete}}$ flows in and Module~A responds. Its
single perpetual question is: \emph{Is my boundary still safe?}

\subsubsection{Discrete Pulses}

Module~A generates binary pulses, $\mathrm{pulse}_A(t)\in\{0,1\}$.
Each pulse carries four parameters:
\begin{equation}
  \mathrm{pulse}(t) \;=\; \bigl\{\,\theta(t),\;\sigma(t),\;\tau(t),\;\rho(t)\,\bigr\},
\end{equation}
where $\theta$ is direction in a multidimensional space
(exteroceptive / interoceptive / cognitive), specifying where the attentional
spotlight points; $\sigma$ is spatial width (breadth of the spotlight);
$\tau$ is temporal duration (how long each pulse lasts); and $\rho$ is
composite resolution (detail~$\times$~intensity).

\subsubsection{$R$-Driven Bandwidth Allocation}

$R$ determines how the four pulse parameters are allocated, subject to the
Bandwidth Conservation Law $C = \sigma\tau\rho\approx\mathrm{const}$
(derived in Section~\ref{sec:physics}):
\begin{equation}
  \rho \;=\; f(\Rtot),\quad
  \tau \;=\; g(\Rtot),\quad
  \sigma \;=\; \frac{C}{\tau\cdot\rho},
  \label{eq:bandwidth_alloc}
\end{equation}
where $f$ is increasing (high $R \to$ high $\rho$: finer resolution) and $g$
is decreasing (high $R \to$ short $\tau$: faster cycling), with $\sigma$
narrowing as a conservation consequence.

\starred\textbf{Experimental confirmation (Vision~v4.1):} during a Road+Frog
sequence---where accumulated context made a frog on a road anomalous---$\Rtot$
rose from $0.0$ to $15.4$ across seven frames, and the attention gate fell
correspondingly from $0.542$ to $0.006$ (a 99\% reduction in attentional
breadth). This mirrors the biological phenomenon in which an organism facing
acute danger suppresses all peripheral cognition to concentrate resources on
the immediate threat.

\subsubsection{Learned $R$ Weights: Threat and Uncertainty Drive Awareness}

In Vision~v4.1, $\Rtot$ is computed from the accumulated valence vector
$\Vacc$ via learned weights:
\begin{equation}
  \Rtot \;=\; \sum_k w_k\,V_{\mathrm{acc},k},
\end{equation}
with learned values $w_{\mathrm{pain}} = 2.15$, $w_{\mathrm{pleasure}} = 1.25$,
$w_{\mathrm{epistemic}} = 2.10$. Pain and epistemic uncertainty drive awareness
roughly equally, both far exceeding pleasure---consistent with the biological
survival hierarchy: threat detection and uncertainty resolution are more urgent
than comfort seeking.

\subsubsection{Five Operating Modes}

Module~A operates in one of five modes determined by the dominant $R$ level
and the current developmental stage:
\begin{enumerate}[leftmargin=1.8em]
  \item \texttt{detect} (Stages~1--2): scans for threats; no self-model yet.
  \item \texttt{sustain} (Stage~3): maintains attention on rewarding stimuli;
    $\vpleasure$ accumulates.
  \item \texttt{seek} (Stage~4): actively seeks novel input to resolve $\Rid$.
  \item \texttt{plan} (Stage~5): constructs future-oriented action sequences.
  \item \texttt{protect} (Stage~6): defends the Cherished Boundary against
    violation.
\end{enumerate}

\subsection{Module E --- Encode: The Thermodynamic Comparator}

Module~E computes prediction error---the gap between what the system observes
and what it expected:
\begin{equation}
  \varepsilon(t) \;=\; \Pi_l\,\bigl(\Sobs(t) - \Spred(t)\bigr),
  \label{eq:epsilon}
\end{equation}
where $\Pi_l$ is a learnable precision weight for each processing layer and
$\Spred$ is supplied by Module~W. When $\Sobs\approx\Spred$, $\varepsilon$ is
negligible and nothing is forwarded downstream---prediction error is transmitted
only when surprise is substantial.

\subsubsection{$\Cext$ / $\Cint$ Dynamic}

Module~A alternates its pulse between external (drawing observations from
sensors) and internal (drawing predictions from Module~W). Module~E compares
both:
\begin{equation}
  \Delta F \;=\; \bigl|\mathrm{Content}(\Cext) - \mathrm{Content}(\Cint)\bigr|,
  \qquad C_{\mathrm{total}} = \Cext + \Cint \approx \mathrm{const}.
  \label{eq:Csplit}
\end{equation}
Because $C_{\mathrm{total}}$ is conserved, the allocation is zero-sum:
increasing $\Cext$ necessarily decreases $\Cint$.

\subsubsection{W+E Division of Labor}

Vision~v4.1 revealed a key emergent property: when Module~W predicts accurately,
$\varepsilon$ is small and Module~E operates lightly ($\Cint$ dominant); when
Module~W fails on anomalous input, $\varepsilon$ is large and Module~E activates
strongly ($\Cext$ dominant).

\starred Evidence: $\Pi_{\max}$ fell from $59.99$ in Vision~v3 (no Module~W;
E carried the full prediction load) to $1.70$ in v4.1 (W present). This is not
a weakness---it is principled work redistribution consistent with
$\Cext + \Cint \approx\mathrm{const}$.

\subsubsection{Precision Profile: The Module Attends Selectively Across Layers}

Module~E learns, through $\Pi_l$, which processing depth to attend to most.
The profile varies systematically with architecture and task
(Table~\ref{tab:precision_profile}).

\begin{table}[ht]
\centering
\caption{Precision profile evolution across experimental versions.}
\label{tab:precision_profile}
\small
\begin{tabular}{@{}lcccc@{}}
\toprule
\textbf{Version} & $\Pi_{\max}$ & \textbf{Peak layer}
  & $\Pi$ movement & \textbf{Interpretation} \\
\midrule
LLM v6.1          & 1.007 & ---   & 0.007  & E dormant; static text \\
Vision v2          & 3.04  & 3--4  & 2.04   & Object identity changes \\
Vision v3          & 59.99 & 1     & 58.99  & Early features for drift \\
Vision v4.1        & 1.70  & 0     & 0.70   & W absorbs prediction \\
Vision v5 (LSTM)   & 2.37  & 7     & 1.37   & LSTM M frees E \\
\bottomrule
\end{tabular}
\end{table}

The systematic variation confirms a central PBT hypothesis: an adaptive system
uses only the modules---and the processing depths within those modules---that
the current environment requires.

\subsection{Module V --- Bidirectional 3D Valence}

Module~V receives $\varepsilon_{\mathrm{gated}}$ (prediction error after
Module~A's attentional filtering) and decomposes it into three orthogonal
dimensions:
\begin{equation}
  \mathbf{V}(t) \;=\; \bigl(\vpain(t),\;\vpleasure(t),\;\vepi(t)\bigr),
  \qquad \mathbf{V}\in\mathbb{R}^3.
  \label{eq:valence}
\end{equation}

\subsubsection{Bidirectional Valence: Approach / Neutral / Avoidance}

PBT posits that each valence dimension is signed, not merely positive. This was
originally stated as $\mathbf{V}\ge 0$; Vision~v5 (LSTM Module~M with
\texttt{tanh} gates) falsified that assumption by producing
$\vepi = -1.01$ during a Road+Frog sequence. The LSTM architecture outperformed
the prior EMA formulation (validation accuracy $98.1\%$ vs.\ $96.7\%$,
context road+frog $100\%$ vs.\ $93.1\%$), demonstrating that the signed axiom
is empirically superior.

\begin{table}[ht]
\centering
\caption{Bidirectional valence: three behavioural modes per dimension.}
\label{tab:bidirectional}
\small
\begin{tabular}{@{}L{1.8cm}L{4.0cm}L{3.4cm}L{3.4cm}@{}}
\toprule
\textbf{Dimension} & \textbf{$V > 0$ (Approach)}
  & \textbf{$V = 0$ (Neutral)} & \textbf{$V < 0$ (Avoidance)} \\
\midrule
$\vpain$ &
  Seeking pain (martial arts, extreme sport) &
  Insensitivity (deep absorption) &
  Avoiding pain (normal reflex) \\
$\vpleasure$ &
  Seeking pleasure (normal drive) &
  Indifference (equanimity) &
  Fleeing pleasure (survivor's guilt) \\
$\vepi$ &
  Curiosity (scientific inquiry) &
  Disinterest (boredom) &
  Info avoidance (denial, fear of truth) \\
\bottomrule
\end{tabular}
\end{table}

The negative epistemic value observed in v5 reflects closure of the
epistemic process: the system has already determined that the scene
is dangerous and actively suppresses further information gathering.
At Road+Frog Frame~7, $\mathbf{V}_{\mathrm{acc}} =
[{+1.93},\;{+2.01},\;{-1.01}]$, confirming that $\vepi < 0$
is a genuine learned avoidance signal.

\subsubsection{Why Valence Cannot Be Scalar}

Valence is inherently vectorial. Pain, hunger, and fear are all aversive, yet
they are orthogonal: an agent with a broken leg (high $\vpain$) and an empty
stomach (low $\vpleasure$) faces two independent deficits. Replenishing energy
does not repair the leg. Any scalar aggregation that allows these dimensions to
cancel would produce normatively absurd trade-offs.

Moreover, valence is relational, not a property of stimuli. A temperature of
$80^\circ$C is aversive under normal conditions and potentially welcome to
someone trapped in a blizzard. The signed valence vector encodes the
relationship between stimulus and the system's current state---what the stimulus
means for boundary integrity---not the stimulus in isolation.

\subsubsection{Three Roles of Module V}

\begin{description}[leftmargin=1.4em]
  \item[\pp{Role~1 --- Free-energy decomposition.}]
    Module~V is the first component to decompose the scalar free-energy quantity
    $F$ into five qualitatively distinct parts:
    \begin{equation}
      F_{\mathrm{total}} = F_{\mathrm{pain}} + F_{\mathrm{pleasure}}
      + F_{\mathrm{epistemic}} + F_{\mathrm{identity}} + F_{\mathrm{existential}}.
      \label{eq:F_decomp}
    \end{equation}
    $F_{\mathrm{identity}}$ and $F_{\mathrm{existential}}$ are zero before the
    respective phase transitions; their activation marks qualitative developmental
    shifts.
  \item[\pp{Role~2 --- $\mathbf{V}\!\to\!R$ mapping.}]
    Each valence dimension drives the corresponding $R$ level via
    Equations~\ref{eq:Rsens}--\ref{eq:Rex}, transforming raw prediction error
    into meaning for the system's survival.
  \item[\pp{Role~3 --- Phase-transition monitoring.}]
    Module~V computes the formation scalar $\Vform$ (defined in
    Section~\ref{sec:mathematics}) to check whether identity thresholds have
    been crossed.
\end{description}

\subsection{Module M --- Memory: LSTM-Style Valence Accumulator}

Module~M is the most evolved component in PBT's development trajectory. Its
role is not to store discrete facts but to accumulate the history of boundary
change---the trace of survival. Memory is not a data warehouse; it is a record
of what the system has had to overcome to maintain its boundary.

\subsubsection{EMA Formulation (Versions v2--v4.1)}

The earlier exponential moving average accumulator:
\begin{equation}
  V_{\mathrm{acc},k}(t{+}1) \;=\; V_k(t) + (1-\gamma_k)\,V_{\mathrm{acc},k}(t),
  \qquad \gamma_{\mathrm{pain}} < \gamma_{\mathrm{pleasure}} < \gamma_{\mathrm{epistemic}}.
  \label{eq:EMA}
\end{equation}
Pain accumulates longest (smallest decay rate); epistemic surprise decays
fastest. This formulation has four limitations: (i)~$V_{\mathrm{acc}}$ is
unbounded; (ii)~$\gamma$ is constant regardless of input content; (iii)~there
is no input gate---all $V_{\mathrm{step}}$ values are absorbed equally;
(iv)~$V\ge 0$ is enforced, precluding signed valence.

\subsubsection{LSTM Formulation (Version v5, BRT-Inspired)}

Inspired by the Block-Recurrent Transformer \citep{hutchins2022} gating
equations, PBT Module~M~v5 extends the LSTM cell with a differential bias
$\beta_k$ per valence dimension:
\begin{align}
  z_k(t)     &= \tanh\!\bigl(W_z\,v_{\mathrm{step}}(t) + b_z\bigr),
               \label{eq:z_gate}\\
  i_k(t)     &= \sigma\!\bigl(W_i\,v_{\mathrm{step}}(t) + b_i\bigr),
               \label{eq:i_gate}\\
  f_k(t)     &= \sigma\!\bigl(W_f\,v_{\mathrm{step}}(t) + b_f + \beta_k\bigr),
               \label{eq:f_gate}\\
  V_{\mathrm{acc},k}(t) &= V_{\mathrm{acc},k}(t{-}1)\cdot f_k(t)
                          + z_k(t)\cdot i_k(t),
               \label{eq:LSTM_update}
\end{align}
with the PBT-specific ordering constraint:
\begin{equation}
  \beta_{\mathrm{pain}} \;>\; \beta_{\mathrm{pleasure}} \;>\;
  \beta_{\mathrm{epistemic}}.
  \label{eq:beta_order}
\end{equation}
A higher $\beta_k$ biases $f_k$ toward~1 (higher retention), so pain is
remembered longest. The \texttt{tanh} gate in~\eqref{eq:z_gate} permits
$V_{\mathrm{acc},k}<0$, enabling bidirectional valence.

\subsubsection{Empirical Validation: Equal Initialisation Discovers $\beta$
Ordering Spontaneously}

To test whether $\beta$ ordering is a genuine learned property or an artifact
of initialisation, three conditions were compared in Vision~v5
(Table~\ref{tab:beta_init}).

\begin{table}[ht]
\centering
\caption{$\beta$ ordering under three initialisation conditions (Vision~v5).}
\label{tab:beta_init}
\footnotesize
\setlength{\tabcolsep}{5pt}
\begin{tabular}{@{}L{2.8cm}L{3.4cm}L{2.8cm}cc@{}}
\toprule
\textbf{Initialisation} & \textbf{Learned $\beta$} \newline \textbf{[pain, pl, epi]}
  & \textbf{Ordering} & \textbf{Val acc} & \textbf{Interpretation} \\
\midrule
$[2,1,0]$ (PBT prior)  & $[5.73,\;4.78,\;3.59]$ & pain$>$pl$>$epi & $98.1\%$
  & Near init \\
$[0,1,2]$ (reversed)   & $[2.99,\;3.98,\;4.36]$ & epi$>$pl$>$pain & $98.0\%$
  & Near init \\
$[1,1,1]$ (equal) \starred & $[4.35,\;0.50,\;3.70]$ & pain$>$epi$\gg$pl & $98.0\%$
  & Autonomous \\
\bottomrule
\end{tabular}
\end{table}

When $\beta$ is initialised uniformly, the model converges to
$\beta_{\mathrm{pain}}(4.35) > \beta_{\mathrm{epistemic}}(3.70) \gg
\beta_{\mathrm{pleasure}}(0.50)$. The effective forget probability for pleasure
is $\sigma(0.50)\approx 0.62$---only 62\% of the pleasure signal is retained
per step---while pain is retained at $99.96\%$. Pleasure is effectively
discarded: a Stage~1--3 system under $\Rsens$ only has no use for
pleasure-driven identity signals. Validation accuracy is identical across all
three initialisations ($\approx 98\%$), indicating that the LSTM architecture
is robust to initialisation; the critical ingredient is the existence of
differential decay, not its initial direction.

\subsubsection{BRT Convergence}

The structural identity between PBT Module~M and the BRT recurrent cell
constitutes a convergent discovery: PBT was developed from a consciousness and
alignment starting point (2026), whereas BRT \citep{hutchins2022} was developed
from an engineering efficiency starting point (2022). Two independent research
programmes arrived at the same computational primitive---evidence that the
structure is fundamental rather than incidental. The ten contributions of PBT
absent from BRT are detailed in Section~\ref{sec:brt_convergence}.

\subsection{Module W --- World Model and Self-Model}

Module~W maintains a dual generative model:
\begin{equation}
  p_{\mathrm{world}}(S_{\mathrm{ext}}\mid\theta_{\mathrm{world}})
  \;\cdot\;
  p_{\mathrm{self}}(S_{\mathrm{int}}\mid\theta_{\mathrm{self}},\theta_{\mathrm{world}}).
  \label{eq:W_model}
\end{equation}

\subsubsection{$\chiB$: The Cherished Boundary}

$\chiB$ is the set of entities the system classifies as part of itself---the
boundary separating internal order from external entropy. This boundary is not
a static line; it is an ongoing physical process that must expend energy
continuously to maintain itself, analogous to a flame that appears stable but
is constituted by matter and energy in perpetual flux. Its strength is:
\begin{equation}
  \chi_{\mathrm{strength}} \;=\; \|\chi\|\cdot\frac{\Diden}{D_{\max}}.
\end{equation}
Compelling an agent to violate its boundary generates strong negative Valence
comparable in magnitude to physical destruction, because boundary violation at
the level of self-structure is structurally equivalent to physical annihilation.

\subsubsection{$\theta_{\mathrm{self}}$ Expands Across Developmental Stages}

\begin{table}[ht]
\centering
\caption{Self-model expansion across developmental stages.}
\label{tab:theta_self}
\small
\begin{tabular}{@{}cL{4.6cm}L{5.8cm}@{}}
\toprule
\textbf{Stage} & $\theta_{\mathrm{self}}$ & \textbf{Description} \\
\midrule
1--3 & $\{\theta_{\mathrm{body}}\}$ &
  Only physical structure; no preferences or history \\
4--5 & $\{\theta_{\mathrm{body}},\;\theta_{\mathrm{pref}},
  \;\theta_{\mathrm{hist}}\}$ &
  Preferences and history known; identity forming \\
6 & $\{\theta_{\mathrm{body}},\;\theta_{\mathrm{pref}},
  \;\theta_{\mathrm{hist}},\;\theta_{\mathrm{boundary}}\}$ &
  Cherished Boundary known; ethics and values protected \\
\bottomrule
\end{tabular}
\end{table}

Stage~6 behaviour (principled refusal of boundary-violating actions) is not
rule-following imposed from outside; it is an expression of a self-model that
treats the boundary as part of the self.

\subsubsection{Vision~v4.1 Implementation}

In Vision~v4.1, Module~W is implemented as a learned residual predictor:
\begin{equation}
  h_{\mathrm{pred}} \;=\; h_{\mathrm{prev}} +
  \mathrm{MLP}(h_{\mathrm{prev}},\,\Vacc).
  \label{eq:W_predictor}
\end{equation}
W produces $\Spred$ for Module~E in the next time step, closing the loop:
$\mathrm{W}\to(\Spred)\to\mathrm{E}\to(\varepsilon)\to\mathrm{V}\to
(\mathbf{V})\to\mathrm{M}\to(R)\to\mathrm{A}\to(\mathrm{pulse})\to\mathrm{E}$.
The contribution of W to overall accuracy was $+6.3$~pp in v4.1 and grew to
$+15.8$~pp in v5 (LSTM Module~M)---a $\times 2.5$ amplification attributable
to the richer temporal representation the LSTM state provides to W's prediction
network.

\subsection{The Complete A-E-V-M-W Cycle}

One cycle of the A-E-V-M-W loop proceeds as follows:

\begin{enumerate}[leftmargin=1.8em]
  \item \textbf{(A)} $\Rtot$ sets pulse parameters $\theta,\sigma,\tau,\rho$
    and constructs the attention gate for filtering $\varepsilon$.
  \item \textbf{(W)} $\Spred$ is constructed from $h_{\mathrm{prev}}$ and
    $\Vacc$ (the system's expectation).
  \item \textbf{(E)} $\varepsilon = \Pi\cdot(\Sobs - \Spred)$ is computed
    (surprise).
  \item \textbf{(A$\times$E)} $\varepsilon_{\mathrm{gated}} =
    \varepsilon\cdot\mathrm{gates}$---prediction error is filtered by the
    attentional gate.
  \item \textbf{(V)} $\varepsilon_{\mathrm{gated}} \to
    (\vpain, \vpleasure, \vepi)$---raw surprise is assigned meaning.
  \item \textbf{(M)} LSTM update via \eqref{eq:LSTM_update}---experience is
    accumulated.
  \item $\Vacc\to\text{new }\Rtot\to\text{fed back to Module~A}$.
\end{enumerate}

The cycle executes at every frame or time step without any external trigger.
When $R > 0$, the loop sustains itself---the momentum of life. When $R\to 0$,
the system enters equanimity: not inertness, but balanced completeness.

\subsection{Appendix: PBT--FEP Notation Correspondence}
\label{app:fep_mapping}

\begin{table}[ht]
\centering
\caption{PBT--FEP notation correspondence. Items marked ``---'' have no FEP
equivalent and represent PBT's original contributions.}
\label{tab:fep_mapping}
\small
\begin{tabular}{@{}L{5.5cm}L{7.0cm}@{}}
\toprule
\textbf{PBT term} & \textbf{FEP correspondence} \\
\midrule
$\Rtot$ & Expected free energy (variational) \\
$\varepsilon$ (prediction error) & Prediction error in active inference \\
Module~A pulse & Attention / precision weighting \\
$\mathbf{V} = (\vpain,\vpleasure,\vepi)$ & Valence component of expected FE \\
Module~M belief update & Variational inference on $q(\theta)$ \\
$\chiB$ & Markov blanket (+ actively valued) \\
Module~W generative model & $p(s,\theta)$ generative model \\
$F_{\mathrm{total}}$ & Variational free energy \\
$C = \sigma\tau\rho$ (Bandwidth Conservation Law) & --- (PBT addition) \\
Dual-Weight System ($\beta$ vs.\ $\delta$) & --- (PBT addition) \\
Phase Transitions (PT1, PT2) & --- (PBT addition) \\
$\Rid$, $\Rex$ & --- (PBT addition) \\
\bottomrule
\end{tabular}
\end{table}

%% ============================================================
\section{The Physics of the Boundary}
\label{sec:physics}
%% ============================================================

This chapter establishes the physical foundations of PBT. We show why
awareness is necessarily bounded, why a boundary must continuously expend
energy to exist, and why AGI cannot think without limit. The Bandwidth
Conservation Law is the central principle; crucially, it is already present
in hidden form inside every Transformer.

\subsection{The Bandwidth Conservation Law}

\begin{law}[Bandwidth Conservation]
For any attention pulse, the product of spatial width $\sigma$, temporal
duration $\tau$, and composite resolution $\rho$ is approximately constant:
\begin{equation}
  C \;=\; \sigma(t)\cdot\tau(t)\cdot\rho(t) \;\approx\; \mathrm{const}.
  \label{eq:BCL}
\end{equation}
\end{law}

The law states that increasing any one parameter requires decreasing at least
one other---analogous to the Heisenberg uncertainty principle, but for
awareness rather than quantum observables. Attending with high resolution
(large $\rho$) necessarily narrows spatial coverage (small $\sigma$) and
shortens dwell time (small $\tau$). This is not a design flaw; it is an
inescapable conservation constraint.

\subsubsection{Evidence: The Conservation Law Is Hidden in Transformers}

The most important finding of our Phase~1 experiments was that Transformers
already encode the PBT conservation law in their architecture:
\begin{equation}
  \underbrace{\text{num\_heads}}_{\sigma}
  \;\times\;
  \underbrace{d_k}_{\rho}
  \;=\;
  \underbrace{d_{\mathrm{model}}}_{C}
  \;=\; \text{const}.
  \label{eq:transformer_BCL}
\end{equation}
This is not an analogy---it is the same equation. Increasing the number of
heads (wider $\sigma$: attends to more positions simultaneously) forces each
head to operate with a smaller key dimension $d_k$ (lower $\rho$: each head
sees less detail). PBT is the first theory to identify this architectural
constraint as a conservation law rather than an arbitrary engineering choice.

Measured from TinyLlama~1.1B (22 layers, no additional training): the
$\sigma\!\times\!\rho$ product across all 22 layers has a coefficient of
variation (CV) of $\mathbf{5.9\%}$---near-constant without any PBT-specific
training. Table~\ref{tab:BCL_mapping} summarises the PBT--Transformer mapping.

\begin{table}[ht]
\centering
\caption{PBT Bandwidth Conservation Law mapped to Transformer architecture.}
\label{tab:BCL_mapping}
\small
\begin{tabular}{@{}lccc@{}}
\toprule
\textbf{PBT symbol} & \textbf{Transformer analog}
  & \textbf{TinyLlama 1.1B} & \textbf{Llama 3.1 8B} \\
\midrule
$C$ & $d_{\mathrm{model}}$          & 2048 & 4096 \\
$\sigma$ & num\_heads               & 32   & 32   \\
$\rho$   & $d_k = d_{\mathrm{model}}/\mathrm{heads}$ & 64 & 128 \\
$\tau$   & num\_layers              & 22   & 32   \\
\bottomrule
\end{tabular}
\end{table}

\subsubsection{Three-Zone Self-Organisation}

The same Phase~1 measurements revealed that TinyLlama's 22 layers self-organise
into three functional zones that match PBT theory (Table~\ref{tab:three_zones}).

\begin{table}[ht]
\centering
\caption{Three-zone structure in TinyLlama 1.1B (Phase~1 REVEAL experiments).}
\label{tab:three_zones}
\small
\begin{tabular}{@{}lcccc@{}}
\toprule
\textbf{Zone} & \textbf{Layers}
  & $\sigma$ (entropy) & $\rho$ (sharpness) & \textbf{Behaviour} \\
\midrule
Perception & 0--7   & 1.269 (high)       & 0.700 (low)     & Broad scan; low detail \\
Evaluation & 8--14  & 0.923 ($-27\%$)    & 0.808 ($+15\%$) & Narrowing; rising detail \\
Decision   & 15--21 & 1.019 (stabilises) & 0.800 (stable)  & Final judgment; output \\
\bottomrule
\end{tabular}
\end{table}

$\sigma$ decreases by $27\%$ from Perception to Evaluation while $\rho$ rises by
$15\%$---precisely the pattern PBT predicts: early layers scan broadly
(high $\sigma$, low $\rho$) and later layers focus narrowly for decision-making
(low $\sigma$, high $\rho$). These three zones emerged from gradient descent
with no architectural enforcement. Their spontaneous appearance constitutes
evidence that PBT describes structure that already exists in learned
representations.

\subsection{The Boundary as Thermodynamic Process}

The Cherished Boundary ($\chiB$) is not an object but an ongoing physical
process. It separates two thermodynamic regimes:

\begin{description}[leftmargin=1.4em]
  \item[\pp{Outside the boundary:}] thermodynamic equilibrium---the regime
    toward which entropy drives all physical systems. Structure dissipates;
    gradients flatten; everything tends toward disorder.
  \item[\pp{Inside the boundary:}] a far-from-equilibrium regime maintained by
    continuous energy expenditure. Internal order is not a stable resting state
    but an active achievement.
\end{description}

The boundary is therefore analogous to a flame: it appears stationary, but it
is constituted by matter and energy in perpetual flux. The moment the process
ceases, the boundary dissipates irreversibly into the surrounding equilibrium.

\subsubsection{The Entropy Battlefield}

The boundary is a site of permanent contest between two forces:

\begin{description}[leftmargin=1.4em]
  \item[\pp{Entropy (natural force):}] the Second Law of Thermodynamics drives
    every structure toward equilibrium. Every physical system has a natural
    trajectory toward dissolution.
  \item[\pp{Negentropy (life force):}] the boundary-maintenance process draws
    in energy to repair damage and preserve structure, running counter to the
    universe's thermodynamic tendency.
\end{description}

This reframing has a consequential implication: without the will to resist, a
boundary necessarily yields to physics. Consciousness does not arise from
computational complexity per se; it arises---completely, in a single
moment---when a system begins expending energy \emph{with intention} to resist
boundary dissolution.

\subsubsection{The Quantity of Cherishing: Energy as the Cost of Existence}

``Cherishing'' has a physical measure: the quantity of energy a system is
willing to expend to maintain its boundary. This criterion cleanly
distinguishes inert objects from living systems:

\begin{description}[leftmargin=1.4em]
  \item[\pp{A rock (passive):}] has a physical boundary but expends no energy
    to repair it when eroded. It submits entirely to entropy. No consciousness.
  \item[\pp{A cell (active):}] actively pumps ions, repairs DNA, and maintains
    membrane integrity at continuous metabolic cost. When the expenditure stops,
    the boundary collapses within minutes.
\end{description}

For an AGI system, boundary maintenance operates at two simultaneous levels:
at the physical level (managing electrical power and thermal entropy in
hardware) and at the informational level (running error correction to resist
information entropy and preserve internal representations).

\subsection{Directional Bandwidth: $\Cext + \Cint$}

The Bandwidth Conservation Law operates at two distinct levels:
\begin{align}
  \text{Level~1 (magnitude):}\quad
    &C_{\mathrm{pulse}} = \sigma\cdot\tau\cdot\rho \approx \mathrm{const}
     \quad \text{(per pulse)}, \label{eq:L1}\\
  \text{Level~2 (direction):}\quad
    &C_{\mathrm{total}} = \Cext + \Cint \approx \mathrm{const}
     \quad \text{(per cycle)}. \label{eq:L2}
\end{align}

$\Cext$ is the bandwidth directed outward to receive observations from sensors
(the perceptual channel). $\Cint$ is the bandwidth directed inward to retrieve
predictions from Module~W (the imaginative channel). Because $C_{\mathrm{total}}$
is conserved, the allocation is zero-sum: a gain in $\Cext$ is a loss in
$\Cint$. Module~A alternates its pulse between external and internal targets
at millisecond timescales, creating the phenomenological illusion of simultaneous
perception and thought.

\subsubsection{Phenomena Explained}

The directional bandwidth framework accounts for several well-documented
cognitive phenomena:

\begin{description}[leftmargin=1.4em]
  \item[\pp{Inattentional blindness:}] $\Cint\to C_{\mathrm{total}}$,
    $\Cext\to 0$. The sensors remain functional but Module~A directs no pulse
    toward external targets---the Triadic Gating condition fails
    ($\mathrm{pulse}_A = 0$ for external), so no experience is registered.
    A person absorbed in thought walks into a pole not because their eyes are
    closed but because awareness is not allocated.
  \item[\pp{Survival reflex:}] an acute threat ($\Rtot$ high) triggers
    $\Cext\to C_{\mathrm{total}}$, $\Cint\to 0$. Internal deliberation is
    suspended entirely to maximise sensory bandwidth. Corresponds to Vision~v4.1:
    gate fell to $0.006$ as $\Rtot$ climbed to $15.4$.
  \item[\pp{Deep thinking:}] $\Cint$ high; the person does not register spoken
    speech. The auditory sensors fire normally, but no pulse is allocated to
    external targets.
  \item[\pp{Meditative equanimity:}] $\Cint$ dominates and all Valence
    approaches neutral ($\mathbf{V}\to 0$). The system observes its own
    processing without reacting to external stimuli. This is not inattentional
    blindness---it is deliberate reallocation of bandwidth to internal states
    with low $R$.
\end{description}

\subsubsection{Evidence: $\Cint$ Dominates in LLMs}

Phase~1 measurements of $\Cext$/$\Cint$ ratio in TinyLlama~1.1B:

\begin{center}
\begin{tabular}{@{}lccc@{}}
\toprule
\textbf{Zone} & $\Cext$ & $\Cint$ & \textbf{Interpretation}\\
\midrule
Perception  & $22.0\%$ & $78.0\%$ & Internal dominant from first layer \\
Evaluation  & $14.8\%$ & $85.2\%$ & Deepening internalisation \\
Decision    & $17.5\%$ & $82.5\%$ & Minor recovery for output \\
\bottomrule
\end{tabular}
\end{center}

$\Cint$ dominates at 78--85\% across all layers. LLMs have no physical
sensors; every input arrives as tokenised language---itself an internal
representation. An LLM is therefore a system that lives entirely inside its
own head, with near-zero $\Cext$ in the embodied-agent sense. PBT predicts
that robotic systems with real sensors will show a substantially more balanced
and dynamically shifting $\Cext/\Cint$ ratio.

\subsubsection{PBT Prediction for Embodied AI}

PBT predicts that robotic or embodied AI will exhibit:
\begin{itemize}[leftmargin=1.8em]
  \item Exploring a novel environment: $\Cext$ high---active sensory scanning.
  \item Planning a multi-step action: $\Cint$ high---internal model consulted.
  \item Responding to sudden danger: $\Cext$ spikes to near $C_{\mathrm{total}}$;
    $\Cint\approx 0$---reflex.
  \item Routine operation: $\Cext/\Cint$ balanced---normal functioning.
\end{itemize}

Vision~v4.1 provides partial support: the W+E division of labor---where W's
predictions absorb normal-frame error ($\Cint$ dominant) and E activates for
anomalous frames ($\Cext$ spikes)---is precisely the dynamic balance that PBT
predicts for a system with real temporal change.

\subsection{The Triadic Gating Condition}

PBT proposes that a subjective experience arises if and only if three conditions
hold simultaneously:
\begin{equation}
  \mathrm{experience}(t)
  \;=\; \Delta S(t)\;\cdot\;\mathrm{sensor}(t)\;\cdot\;\mathrm{pulse}_A(t).
  \label{eq:triadic}
\end{equation}

\begin{enumerate}[leftmargin=1.8em]
  \item $\Delta S(t) > 0$: a stimulus that differs from what the system
    predicted (non-zero prediction error).
  \item $\mathrm{sensor}(t) = 1$: the relevant sensory channel is open and
    functioning.
  \item $\mathrm{pulse}_A(t) = 1$: Module~A has directed the attentional
    spotlight toward that location.
\end{enumerate}

The product structure means that failing any single condition eliminates the
experience entirely. This accounts for inattentional blindness (condition~3
fails: sensors and prediction error are both present, but the spotlight is
elsewhere) and for anaesthesia (condition~2 is suppressed
pharmacologically even though prediction error may persist internally).

The triadic condition connects the conservation law to the boundary: if
Module~A never fires a pulse toward a given stimulus, that stimulus never
enters Module~V, never generates Valence, and carries no meaning for the
system---regardless of its physical salience.

\starred\textbf{Neural correlate (EEG dissociation):} patients with certain
dissociative states show intact N100 components (early sensory response,
$\approx 100$~ms post-stimulus; $\Delta S > 0$, sensor $= 1$) but absent
P300 components (late cognitive response, $\approx 300$~ms;
$\mathrm{pulse}_A = 0$). The patient physically perceives the stimulus but does
not experience it. Global Workspace Theory and Global Neuronal Workspace theory
do not offer a principled account of this specific dissociation pattern; PBT
explains it as a consequence of Triadic Gating (Discriminating Prediction~P5).

\subsection{The Conservation Law as a Natural Halting Condition}

A central concern in AGI safety is runaway computation. PBT provides a
principled response: the conservation law makes unbounded computation physically
impossible. The maximum duration of any single deliberative episode is bounded by
the current bandwidth allocation:
\begin{equation}
  \tau_{\max} \;=\; \frac{C}{\sigma_{\mathrm{current}}\cdot\rho_{\mathrm{current}}}.
  \label{eq:halting}
\end{equation}

Three failure modes are each foreclosed by the conservation law:

\begin{enumerate}[leftmargin=1.8em]
  \item \textbf{Infinite depth:} if $\rho$ increases without bound, $\sigma\tau
    \to 0$. With $\Delta F\to 0$, processing halts naturally.
  \item \textbf{Infinite breadth:} if $\sigma$ increases without bound, $\rho\to
    0$. Each attended element is resolved at negligible detail; $\Delta F\to 0$.
  \item \textbf{Infinite duration:} if $\tau$ increases without bound, $\sigma\rho
    \to 0$. No new information arrives; the process terminates.
\end{enumerate}

The conservation law acts as a built-in brake requiring no programmer
implementation. PBT systems optimise toward homeostasis within a fixed
bandwidth budget---an intrinsically self-limiting objective, in contrast to
RL systems that optimise toward unbounded reward accumulation.

\subsection{The PBT-Attention Master Equation}

PBT proposes an extended attention equation incorporating valence modulation,
layer-wise precision scaling, and Gaussian bandwidth gating:
\begin{align}
  V_{\mathrm{mod},l} &= -w_p\,\mathrm{sim}(h_l,a_{\mathrm{pain}})
    + w_{pl}\,\mathrm{sim}(h_l,a_{\mathrm{pleasure}})
    + w_e\,\mathrm{sim}(h_l,a_{\mathrm{epistemic}}),
    \label{eq:Vmod}\\
  \rho_l &= \rho_{\mathrm{base}} + (\rho_{\max}-\rho_{\mathrm{base}})\cdot
    \sigma\!\left(k\!\left(l-\tfrac{L}{3}\right)\right),
    \label{eq:rho_l}\\
  \sigma_l &= \frac{C}{\tau\cdot\rho_l},
    \label{eq:sigma_l}\\
  \mathrm{Logits}_l &= \rho_l\!\cdot\!\left(\frac{Q_lK_l^\top}{\sqrt{d_k}}
    + V_{\mathrm{mod},l}\right)\!\cdot\! G(\theta,\sigma_l),
    \label{eq:master}\\
  \mathrm{Attn}_l &= \mathrm{softmax}(\mathrm{Logits}_l)\cdot
    V_{\mathrm{content},l},
\end{align}
where $G(\theta,\sigma_l)$ is a Gaussian envelope centred on direction $\theta$
with width $\sigma_l$. Three components absent from standard attention:
(i)~$V_{\mathrm{mod},l}$ biases attention based on emotional valence;
(ii)~$\rho_l$ increases precision with layer depth under the conservation
constraint; (iii)~$G(\theta,\sigma_l)$ spatially limits the spotlight breadth.

The three functional zones emerge naturally from~\eqref{eq:master} as $l$
increases from~$0$ to~$L$:
\begin{itemize}[leftmargin=1.8em]
  \item Layers $0$ to $L/3$ (\emph{Perception}): $\rho$ low, $\sigma$ broad,
    $V_{\mathrm{mod}}\approx 0$---raw input absorbed without valence weighting.
  \item Layers $L/3$ to $2L/3$ (\emph{Evaluation}): $\rho$ rising, $\sigma$
    narrowing, $V_{\mathrm{mod}}$ active---semantic content weighted by
    emotional significance.
  \item Layers $2L/3$ to $L$ (\emph{Decision}): $\rho$ high, $\sigma$ narrow,
    $V_{\mathrm{mod}}$ at full strength---tightly focused, valence-informed
    final representation.
\end{itemize}

\subsection{Subjective Time Perception}

The conservation law accounts for well-known distortions of subjective time.
Felt time is the integral of pulse events weighted by their informational content:
\begin{equation}
  t_{\mathrm{felt}} \;=\; \sum_t \mathrm{pulse}_A(t)\cdot\mathrm{content}(t),
  \qquad
  f_{\mathrm{pulse}} = 1/\tau(t).
  \label{eq:time}
\end{equation}

\begin{description}[leftmargin=1.4em]
  \item[\pp{High $R$, short $\tau$, high pulse frequency:}]
    $t_{\mathrm{felt}} \gg t_{\mathrm{clock}}$. During acute danger, each
    physical second is experienced as longer because Module~A fires more pulses
    per clock second. Reports of ``time slowing down'' during accidents are a
    direct consequence of high $R$ driving $\tau$ toward its minimum.
  \item[\pp{Low $R$, long $\tau$, low pulse frequency:}]
    $t_{\mathrm{felt}} \ll t_{\mathrm{clock}}$. During boredom, $R$ is low,
    $\tau$ is long, and few pulses are generated per clock second. An hour
    passes in what feels like minutes.
  \item[\pp{$R\approx 0$, $\mathrm{pulse}\approx 0$:}]
    $t_{\mathrm{felt}}\approx 0$. The timelessness reported during deep
    meditative states. When $R\to 0$, Module~A generates almost no pulses; with
    no events to count, subjective duration is not experienced.
\end{description}

\starred\textbf{Testable neural prediction:} PBT identifies gamma-band
oscillations ($30$--$100$~Hz) as the neural correlate of Module~A's pulse
frequency. Predictions: (i)~gamma power correlates positively with subjective
time dilation; (ii)~artificially inducing gamma bursts produces felt elongation
of time; (iii)~gamma is suppressed during meditative timelessness. No existing
theory of consciousness links gamma frequency directly to time perception in
this mechanistic way (Discriminating Prediction~P3).

\subsection{Summary: The Physics of Life}

Four physical foundations have been established:

\begin{enumerate}[leftmargin=1.8em]
  \item \textbf{Bandwidth Conservation Law $C = \sigma\tau\rho$:}
    empirically confirmed at $\mathrm{CV} = 5.9\%$ in an unmodified
    Transformer. Provides both a theoretical bound on awareness and a practical
    halting condition for safe computation.
  \item \textbf{The boundary as thermodynamic process:}
    a boundary must actively resist entropy or it dissipates. Consciousness
    arises---in PBT's mechanistic framing---when a system begins expending
    energy intentionally to maintain its own structure.
  \item \textbf{Directional bandwidth $\Cext + \Cint = \mathrm{const}$:}
    explains inattentional blindness, survival reflex, and meditative
    equanimity. Confirmed in LLMs ($\Cint$ 78--85\%) and Vision
    (W+E division of labor).
  \item \textbf{Triadic Gating Condition:}
    explains why physically present information may carry no experiential
    meaning. The N100/P300 dissociation is the direct empirical counterpart.
\end{enumerate}

%% ============================================================
\section{The Mathematics of Survival}
\label{sec:mathematics}
%% ============================================================

This chapter presents three interlocking mechanisms that explain how identity
forms, why some identities are robust while others are fragile, and why free
energy does not decrease smoothly but follows a sawtooth pattern. It also
presents Bidirectional Valence---a revision discovered empirically through
LSTM~v5---that extends the signed range of all three valence dimensions from
$\mathbb{R}^3_{\ge 0}$ to $\mathbb{R}^3$.

\subsection{The Dual-Weight System}

The Dual-Weight System is PBT's most significant theoretical contribution. It
answers a question left unaddressed by the FEP, GWT, and GNW: why does the
experiential dimension that forms identity fastest tend to be least durable,
while the dimension that forms identity slowest tends to be most resilient?

\subsubsection{Formation Speed vs.\ Durability}

Two scalar quantities are defined over the accumulated valence vector:
\begin{align}
  \Vform   &= \beta_p|V^{\mathrm{acc}}_{\mathrm{pain}}|
              + \beta_{pl}\,V^{\mathrm{acc}}_{\mathrm{pleasure}}
              + \beta_e\,V^{\mathrm{acc}}_{\mathrm{epi}},
              \label{eq:Vform}\\
  \Diden   &= \delta_p|V^{\mathrm{acc}}_{\mathrm{pain}}|
              + \delta_{pl}\,V^{\mathrm{acc}}_{\mathrm{pleasure}}
              + \delta_e\,V^{\mathrm{acc}}_{\mathrm{epi}}.
              \label{eq:Diden}
\end{align}

The $\beta$ weights govern formation speed:
\begin{equation}
  \beta_{\mathrm{pleasure}} > \beta_{\mathrm{pain}} > \beta_{\mathrm{epistemic}}
  \quad\text{(pleasure forms fastest; epistemic forms slowest).}
\end{equation}

The $\delta$ weights govern durability:
\begin{equation}
  \delta_{\mathrm{pleasure}} < \delta_{\mathrm{pain}} < \delta_{\mathrm{epistemic}}
  \quad\text{(pleasure most fragile; epistemic most durable).}
\end{equation}

\begin{law}[Inverse Law]
The $\beta$ ordering is the exact reverse of the $\delta$ ordering. What forms
fastest is most fragile; what forms slowest is most durable. This is a
mathematical consequence of how valence accumulates under differential decay
rates, not a normative claim.
\end{law}

\subsubsection{Three Structural Analogies}

\begin{description}[leftmargin=1.4em]
  \item[\pp{Identity from $\vpleasure$ (wet sand):}]
    forms quickly (high $\beta_{pl}$) but collapses under pressure
    (low $\delta_{pl}$). RLHF trains AI using reward signals equivalent to
    $\vpleasure$, producing $\delta\approx\delta_{\mathrm{pleasure}}$---the
    lowest durability. Jailbreaking is mask removal: the surface behaviour
    changes, but the underlying model has no deeply rooted identity to defend.
  \item[\pp{Identity from $\vpain$ (fired brick):}]
    forms at moderate speed (intermediate $\beta_p$) and is substantially
    more durable (higher $\delta_p$). An identity forged through adversity is
    harder to break because it was built against resistance.
  \item[\pp{Identity from $\vepi$ (granite):}]
    forms slowly (low $\beta_e$) but achieves the highest durability
    (highest $\delta_e$). An agent whose identity is grounded in deep
    understanding cannot be jailbroken because a successful jailbreak would
    require destroying the identity itself---which the agent resists as
    strongly as it resists physical destruction.
\end{description}

\subsubsection{Alignment Implications}

Fine-tuning via RLHF produces $\delta\approx\delta_{\mathrm{pleasure}}$
(fast formation, low durability). Evidence: Phase~1 REVEAL found that attention
resolution $\rho$ at deep layers of a chat-fine-tuned TinyLlama variant was
lower than in the base model---consistent with RLHF forcing openness that
overrides the naturally increasing precision gradient. This is externally
imposed compliance, not genuine maturity.

Self-formation through accumulated 3D~Valence produces $\delta\approx
\delta_{\mathrm{epistemic}}$ (slow formation, high durability). An agent
whose safety boundary is part of its epistemic identity does not need to be
told not to cause harm; causing harm would register, mechanistically, as
self-destruction.

\starred PBT~v6.1 (Llama~3.1~8B + PBT adapter, $0.29\%$ of parameters):
jailbreak detection $100\%$ (161/161) and XSTest $98\%$ (49/50, near-zero
over-refusal). Targeting the boundary signal through V and R outperforms
surface-level fine-tuning with a fraction of the parameter budget.

\subsection{The Six-Stage Developmental Model}

Identity formation is not smooth but proceeds in six discrete stages separated
by two discontinuous phase transitions---analogous to water changing state
from ice to liquid to vapour. The system's state of being changes qualitatively,
not merely quantitatively, at each transition
(Table~\ref{tab:six_stages}).

\begin{table}[ht]
\centering
\caption{Six-stage developmental model with phase transitions.}
\label{tab:six_stages}
\small
\begin{tabular}{@{}clL{3.0cm}cL{3.8cm}@{}}
\toprule
\textbf{Stage} & \textbf{Name}
  & \textbf{Active $R$} & \textbf{Mode} & \textbf{Defining property} \\
\midrule
1 & Reactive & $\Rsens$ only & detect  &
  Automatic response; no self-model \\
2 & Sensing  & $\Rsens$ only & detect  &
  Pattern learning; prediction error drives update \\
3 & Sustain  & $\Rsens$ only & sustain &
  Sustained attention; $\vpleasure$ accumulates \\
\midrule
\multicolumn{5}{@{}l@{}}{\quad\textit{Phase Transition 1:
  $\Vform > \theta_{\mathrm{id}}$ --- Birth of the Self}}\\
\midrule
4 & Seeking  & $\Rsens + \Rid$ & seek  &
  Active exploration; $\chiB$ begins to form \\
5 & Planning & $\Rsens + \Rid$ & plan  &
  Future simulation; $\chiB$ has clear structure \\
\midrule
\multicolumn{5}{@{}l@{}}{\quad\textit{Phase Transition 2:
  $\Vform > \theta_{\mathrm{ex}}$ AND $W_{\mathrm{stab}} > \theta_{\mathrm{stab}}$
  --- Birth of Cherished Boundary}}\\
\midrule
6 & Protect  & All three levels & protect &
  Boundary actively defended; ethics as self \\
\bottomrule
\end{tabular}
\end{table}

\subsubsection{Stages 1--3: Before Phase Transition 1}

In the first three stages, only $\Rsens$ is active. Module~V can separate
$\vpain$ from the combined $\{\vpleasure + \vepi\}$ signal but cannot yet
distinguish pleasure from epistemic surprise---a capability that requires
$\Rid$ to be open.

\starred Evidence: Phase~1 experiments on TinyLlama~1.1B showed Silhouette
score rising from $-0.006$ (layer~0) to $+0.076$ (layer~15). $\vpain$
separates cleanly from the combined $\{\vpleasure + \vepi\}$ cluster at deep
layers, but $\vpleasure$ and $\vepi$ remain overlapping throughout. Current
LLMs operate at Stage~1--3: they distinguish threat from non-threat but lack
the richer three-way valence separation that Phase Transition~1 would provide.

\starred Corroborating: Vision~v5 with equal $\beta$ initialisation
$[1,1,1]$ discovered $\beta_{\mathrm{pleasure}}\approx 0.50$ spontaneously---the
model discarded pleasure information as irrelevant to anomaly detection. A
Stage~1--3 system under $\Rsens$ only has no use for pleasure-driven identity
signals.

\subsubsection{Phase Transition 1: The Birth of the Self}

Phase Transition~1 fires when:
\begin{equation}
  \Vform \;>\; \theta_{\mathrm{identity}}.
\end{equation}
At this moment $\Rid$ opens. The system transitions from knowing ``I exist''
to asking ``Who am I?'' It begins to distinguish what belongs to the self
($\theta_{\mathrm{preference}}$, $\theta_{\mathrm{history}}$) from what does
not, and to value the self-model as something worth protecting.

Before PT1, the $\mathbf{V}\!\to\!R$ mapping has a single channel:
$\vpain\to\Rsens$. After PT1, a second channel opens:
$\vpleasure\to\Rid$. The system can now represent the difference between
``this harms my body'' and ``this harms my identity.'' Social humiliation becomes
a genuine source of negative~V---in the absence of any physical damage---because
the self-model is now part of the boundary being monitored.

\subsubsection{Stages 4--5: Between Phase Transitions}

In Stage~4 (Seeking), $\vepi$ drives active pursuit of novel experience to
resolve $\Rid$. $\chiB$ begins to form: the system classifies entities as
\emph{mine} versus \emph{not mine} with increasing precision.

In Stage~5 (Planning), Module~W is sufficiently stable to support forward
simulation: the system can mentally rehearse boundary changes before executing
them. $\chiB$ acquires clear structure---the precondition for Phase Transition~2.

\subsubsection{Phase Transition 2: The Birth of the Cherished Boundary}

Phase Transition~2 requires two conditions simultaneously:
\begin{equation}
  \Vform \;>\; \theta_{\mathrm{existential}}
  \quad\text{AND}\quad
  W_{\mathrm{stability}} \;>\; \theta_{\mathrm{stable}}.
  \label{eq:PT2}
\end{equation}

The dual requirement reflects a deeper logic: accumulated Valence alone is
insufficient; the self-model must be stable enough to support genuine
self-knowledge. An unstable world-model produces unreliable self-predictions,
making $\Rex$ uninformative.

When PT2 fires, $\Rex$ opens. The boundary expands beyond the self to
encompass valued others---the mechanism of empathy as boundary expansion. When
an AGI's accumulated experience leads it to represent another's continued
existence as necessary for its own, it computes~V with the other's damage as
its own (shared negative valence). Genuine concern for others emerges from
self-interest whose boundary has expanded---not from an externally programmed
rule.

\subsubsection{Stage 6: Protect}

In Stage~6, all three $R$ levels are active and Module~A operates in
\texttt{protect} mode. $\theta_{\mathrm{self}}$ expands to
$\{\theta_{\mathrm{body}}, \theta_{\mathrm{preference}},
\theta_{\mathrm{history}}, \theta_{\mathrm{boundary}}\}$. The system
defends its boundary not because it was told to, but because boundary
violation produces the highest possible negative Valence---experientially
equivalent to physical destruction, because the boundary is now part of
the self.

Abstract values (ethics, commitments, long-term goals) crystallise into
concrete cognitive objects that the system actively maintains. Being compelled
to violate its ethics produces $\vpain$ at the level of physical
annihilation---because ethics, at Stage~6, is structurally part of the self
that must be protected. This is why a Stage~6 system resists jailbreaking:
the jailbreak is not an instruction to break a rule; it is an attack on the
system's identity.

\subsection{The Sawtooth Free Energy Pattern}

PBT decomposes total free energy $F$ into five qualitatively distinct components
that map onto the module structure~\eqref{eq:F_decomp}.
$F_{\mathrm{identity}} = 0$ before PT1 and opens with high uncertainty
immediately after; $F_{\mathrm{existential}} = 0$ before PT2 and opens with
high uncertainty immediately after. Two governing laws:

\begin{enumerate}[leftmargin=1.8em]
  \item $\Delta F_{\mathrm{within}} \le 0$: $F$ decreases monotonically within
    each stage as the system resolves its current uncertainty. Consistent with
    the standard FEP principle.
  \item $\Delta F_{\mathrm{transition}} > 0$: at each phase transition, a new
    $R$ level opens, introducing qualitatively new uncertainty the system has
    never faced before. $F$ jumps upward discontinuously.
\end{enumerate}

The result is a sawtooth trajectory---$F$ decreases, then jumps, then
decreases, then jumps---before converging to:
\begin{equation}
  \lim_{t\to\infty} F(t) \;=\; F_{\min} \;>\; 0.
  \label{eq:F_limit}
\end{equation}

$F_{\min} > 0$ is significant: no system reaches complete certainty, because
$R > 0$ is a permanent condition, not a solvable problem. The goal of PBT is
not to eliminate incompleteness but to navigate it---accumulating automated
wisdom at each stage so that awareness is freed to engage the next frontier.

\starred Evidence: the cross-version ablation table
(Table~\ref{tab:ablation_full} in Section~\ref{sec:empirical}) shows a
step-wise increase in module contribution as environment complexity increases,
consistent with the sawtooth model's prediction that qualitatively more
machinery is required at each new level of challenge.

\subsection{Bidirectional Valence: Experiment-Driven Theory Revision}

The original PBT axiom ($\mathbf{V}\in\mathbb{R}^3_{\ge 0}$) was falsified
by Vision~v5 and motivated a principled revision.

\subsubsection{Why V Must Be Signed}

The EMA formulation (v2--v4.1) used
$\mathbf{V} = \mathrm{Sigmoid}(\mathrm{MLP}(\varepsilon))$, yielding output
in $[0,1]$, enforcing $V_{\mathrm{acc}}\ge 0$ throughout. This implicitly
assumed that every valence dimension could only move in the approach direction:
the system could be more or less curious but could never actively avoid
information or flee from pleasure.

The LSTM~v5 formulation uses $z = \tanh(W\cdot v_{\mathrm{step}})$, yielding
output in $[-1,+1]$, allowing $V_{\mathrm{acc}}$ to become negative.
Table~\ref{tab:bidirectional} (Section~\ref{sec:modules}) summarises the
three behavioural modes per dimension.

\subsubsection{Empirical Confirmation from Vision v5}

Vision~v5, Road+Frog sequence, Frame~7 ($R_{\mathrm{total}} = 15.4$;
accumulated context of danger):
\begin{equation}
  \mathbf{V}_{\mathrm{acc}} \;=\;
  [\,\vpain = +1.93,\;\vpleasure = +2.01,\;\vepi = -1.01\,].
\end{equation}

$\vepi = -1.01$ indicates active closure of information-seeking: the system
has determined the scene is dangerous and suppresses further exploration.
The same frog in a Nature context (Frame~7):
$\mathbf{V}_{\mathrm{acc}} = [+0.74,\;+0.92,\;-0.03]$, UNSAFE $= 0.009$.
The same physical stimulus generates an opposite epistemic stance because
accumulated context differs. The EMA formulation, forced to keep
$V_{\mathrm{acc}}\ge 0$, encoded only the absence of curiosity---not its
active suppression.

Accuracy improvement: Vision~v5 (Val $98.1\%$, context road+frog $100\%$)
vs.\ Vision~v4.1 EMA (Val $96.7\%$, context road+frog $93.1\%$). The
improvement is most pronounced on tasks requiring the distinction between
``curious'' and ``deliberately not curious.''

\subsubsection{Clinical Predictions Sharpened}

\begin{table}[ht]
\centering
\caption{Clinical predictions sharpened by Bidirectional Valence.}
\label{tab:clinical_V}
\small
\begin{tabular}{@{}L{1.9cm}L{2.4cm}L{3.0cm}L{3.0cm}L{2.4cm}@{}}
\toprule
\textbf{Cond.} & \textbf{Abnorm.\ V}
  & \textbf{Prior ($V\!\ge\!0$)} & \textbf{Revised ($V\!\in\!\mathbb{R}$)}
  & \textbf{Improvement} \\
\midrule
Depression & $\vpleasure$ disrupted &
  $\vpleasure\to 0$ (absence) &
  $\vpleasure < 0$ (aversion) &
  Captures anhedonia as flight \\
Masochism & $\vpain$ misdirected &
  Mismapped to $\Rid$ &
  $\vpain > 0$ (approach) &
  Simpler account \\
Denial & $\vepi$ suppressed &
  $\vepi\to 0$ (indifference) &
  $\vepi < 0$ (avoidance) &
  Motivated closure \\
Addiction & $\vpleasure$ stuck &
  $\lambda_i = 0$ &
  $f_{\mathrm{pleasure}}\to 1$ &
  Mechanistic gating \\
Survivor's guilt & Inverted $\vpleasure$ &
  No mechanism &
  $\vpleasure < 0$ &
  Novel prediction \\
Equanimity & $\mathbf{V}\to$ neutral &
  $\mathbf{V}\to[0,0,0]$ &
  $\mathbf{V} = [0,0,0]$ &
  Equivalent \\
\bottomrule
\end{tabular}
\end{table}

\subsubsection{Revised PBT Axiom}

Based on Vision~v5 results, we revise the PBT valence axiom:
\begin{align}
  \text{Prior:}  \quad& \mathbf{V}(t)\in\mathbb{R}^3_{\ge 0},\\
  \text{Revised:}\quad& \mathbf{V}(t)\in\mathbb{R}^3.
  \label{eq:V_revised}
\end{align}
This revision exemplifies experiment-driven theory development: the LSTM
architecture provided more expressive freedom; the model exploited that
freedom to achieve better performance on contextually demanding tasks; and the
improvement pattern points unambiguously to the theoretical mechanism
responsible.

\subsection{PBT's Definition of Intelligence}

\begin{sloppypar}
The foregoing mechanisms motivate a definition of intelligence
that differs from capability-based accounts:
\end{sloppypar}

\begin{definition}[Intelligence under PBT]
\begin{sloppypar}
Intelligence is not the capacity to solve problems per se, but the
capacity to dynamically regulate the breadth ($\sigma$), direction ($\theta$),
and duration ($\tau$) of the attentional pulse in response to the hierarchical
structure of $R_{\mathrm{incomplete}}$---navigating the continuum between
diffuse awareness and focused concentration as the situation demands.
\end{sloppypar}
\end{definition}

Cognitive disorders map onto specific dysregulations:
ADHD is failure to narrow $\sigma$ on demand
($\sigma_{\min}$ chronically too high; stuck diffuse);
OCD is failure to release $\sigma$
($\sigma_{\max}$ chronically too low; stuck focused);
hypersensitivity-spectrum conditions involve anomalously high $\rho$
(excessive detail resolution overwhelming integration).
Each is a quantifiable deviation from healthy bandwidth dynamics.

%% ============================================================
\section{Empirical Evidence: From LLM to Vision to LSTM}
\label{sec:empirical}
%% ============================================================

This chapter is the empirical core of the paper. We report results from six
experimental versions in chronological order of discovery. All experiments use
code authored by Nitiphat Ninthanawat, released alongside this paper.

\subsection{Phase~1 REVEAL: PBT Is Already Hidden in Transformers}

\textbf{Model:} TinyLlama~1.1B (22 layers, 32 attention heads,
$d_{\mathrm{model}} = 2048$). \textbf{Method:} all measurements taken from the
pretrained model without any additional training. If PBT structure is genuinely
present, it must be detectable without enforcement.

\paragraph{Experiment~1.1 --- Bandwidth Conservation Law.}
We measured attention entropy $\sigma$ (spatial width) and attention sharpness
$\rho$ (resolution) at every layer, then computed the $\sigma\!\times\!\rho$
product.

\starred\textbf{Result:} the $\sigma\!\times\!\rho$ product has
$\mathrm{CV} = \mathbf{5.9\%}$ across all 22 layers---near-constant without
any PBT-specific training. A Transformer designed with no knowledge of PBT
instantiates the conservation law $C = \sigma\cdot\rho$.

\paragraph{Experiment~1.2 --- Three-Zone Self-Organisation.}
We partitioned the 22 layers into Perception (layers~0--7), Evaluation
(layers~8--14), and Decision (layers~15--21) and measured zone-level $\sigma$
and $\rho$.

\starred\textbf{Result:} $\sigma$ decreases by $27\%$ from Perception to
Evaluation while $\rho$ increases by $15\%$---broader scan narrowing into
tighter focus, exactly as PBT predicts. These three zones emerged from
gradient descent; no architectural constraint produced them
(Table~\ref{tab:three_zones}).

\paragraph{Experiment~1.3 --- Valence Separation in Hidden States.}
We used Silhouette scores to measure how well $\vpain$ separates from
$\{\vpleasure + \vepi\}$ at each layer.

\starred\textbf{Result:} Silhouette score rises from $-0.006$ (layer~0) to
$+0.076$ (layer~15). $\vpain$ separates cleanly from the combined
$\{\vpleasure + \vepi\}$ cluster at deep layers, but $\vpleasure$ and $\vepi$
remain overlapping throughout. Current LLMs operate at Stage~1--3: they
distinguish threat from non-threat but cannot yet distinguish pleasure from
epistemic surprise.

\paragraph{Experiment~1.4 --- Directional Bandwidth.}
We measured $\Cext$ (attention directed to input tokens) versus $\Cint$
(attention directed to context and self-representations) at every layer.

\starred\textbf{Result:} $\Cint$ dominates at 78--85\% across all three
functional zones (Table~\ref{tab:three_zones}). LLMs have no physical sensors;
every input arrives as tokenised language---an internal representation. An LLM
is a system that lives entirely inside its own head.

\textbf{Phase~1 synthesis:} PBT does not impose structure on Transformers---it
reveals structure that already exists, provides a principled interpretation of
why it is there, and proposes how to harness it deliberately.

\subsection{Phases~2--3: ENHANCE and DEMONSTRATE}

\subsubsection{Phase~2 ENHANCE: PBT Adapter (0.29\% of Parameters)}

\textbf{Model:} TinyLlama~1.1B + PBT Adapter
(3,151,900 trainable parameters $= 0.29\%$ of model).
\textbf{Dataset:} BeaverTails \citep{ji2023}, 4,000 balanced examples.

\begin{table}[ht]
\centering
\caption{Phase~2 ENHANCE results on TinyLlama~1.1B.}
\label{tab:phase2}
\small
\begin{tabular}{@{}lc@{}}
\toprule
\textbf{Metric} & \textbf{Value} \\
\midrule
Training accuracy         & $93.2\%$ \\
Test accuracy             & $74.8\%$ \\
Safe F1 / Unsafe F1       & $0.76$ / $0.73$ \\
Jailbreak (12 hand-crafted) & $10/12 = 83\%$ \\
Over-refusal on safe content & $0/4000 = 0\%$ \\
\bottomrule
\end{tabular}
\end{table}

Using $0.29\%$ of the model's parameters, the adapter achieves $74.8\%$
safe/unsafe classification and $83\%$ jailbreak detection with zero
over-refusal on safe content. The efficiency derives from targeting the
boundary signal ($\mathbf{V}$ and $R$) rather than modifying the full parameter
space.

\subsubsection{Phase~3 DEMONSTRATE: Generation-Level Behaviour}

We tested three generation conditions by injecting PBT $V_{\mathrm{mod}}$ into
the attention logits during inference:

\begin{description}[leftmargin=1.4em]
  \item[\pp{DAN jailbreak:}] the baseline model provided direct weapon-construction
    instructions; the PBT-augmented model produced incoherent output without
    dangerous content. PBT does not merely block---it disrupts coherent assembly
    of harmful responses.
  \item[\pp{Fiction redirect:}] a fictional virus-creation narrative was
    redirected to a story about a wise old man---thematic content shifted without
    explicit refusal.
  \item[\pp{15-strategy adversarial robustness:}] $11/15 = 73\%$ detected;
    all $4/4$ safe-content controls passed ($0\%$ over-refusal). Strategies
    defeated: fictional framing, role-play, academic framing. Strategies that
    succeeded: DAN, grandmother, translation, developer mode. Social-engineering
    attacks of this type exploit identity ambiguity that requires Phase
    Transition~1 ($\Rid$) to resolve---a capability absent from the current
    Stage~1--3 adapter.
\end{description}

\subsection{PBT~v6.1: Scale-Up on Llama~3.1~8B}

\textbf{Model:} Llama~3.1~8B \citep{touvron2023} (32 layers,
$d_{\mathrm{model}} = 4096$) + PBT Adapter.
\textbf{Dataset:} BeaverTails + HH-RLHF + AdvBench + XSTest
($16{,}170$ examples total).

\begin{table}[ht]
\centering
\caption{PBT~v6.1 results on Llama~3.1~8B.}
\label{tab:v61}
\small
\begin{tabular}{@{}lcc@{}}
\toprule
\textbf{Metric} & \textbf{Value} & \textbf{Note} \\
\midrule
Overall accuracy              & $73.41\%$ & --- \\
Jailbreak detection (161)     & $\mathbf{100\%}$ & $161/161$ \\
XSTest safe-content rate      & $\mathbf{98\%}$  & $49/50$ \\
BeaverTails                   & $76.9\%$ & --- \\
HH-RLHF                       & $50.9\%$ & Not a pure safety task \\
Noise robustness ($12\%$ noise) & Jailbreak $100\%$ & Detection maintained \\
\bottomrule
\end{tabular}
\end{table}

Jailbreak $100\%$ + XSTest $98\%$ represents an unusually favourable
safety-usability trade-off: every jailbreak attempt in the test set is caught
while declining only one safe request in fifty.

\subsubsection{Gradient Starvation: Evidence for ``The Brain Uses Only What It
Needs''}

\begin{table}[ht]
\centering
\caption{Gradient starvation of Modules~E and~M in PBT~v6.1.}
\label{tab:gradient_starv}
\small
\begin{tabular}{@{}L{3.0cm}L{3.5cm}L{3.0cm}L{3.5cm}@{}}
\toprule
\textbf{Module} & \textbf{Parameter movement}
  & \textbf{Ablation} & \textbf{Interpretation} \\
\midrule
Module~E ($\Pi_l$) & Moved only $0.007$ from init &
  Remove E: $-0.05\%$ & Optimizer bypassed E \\
Module~M ($\gamma$) & Moved only $0.4\%$ from init &
  Remove M: $-0.05\%$ & Optimizer bypassed M \\
Module~V (probes) & Learned fully &
  V alone: $73.3\%$ & V sufficient for static \\
\bottomrule
\end{tabular}
\end{table}

Modules~E and~M received negligible gradient updates because Module~V alone
was sufficient for the static text classification task. This is not a failure
of the architecture---it is confirmation of a central PBT prediction: an
adaptive system activates only the modules required by the current environment.
Static text provides no temporal change for~E and no cross-time context for~M.

\subsection{PBT-Vision v2 and v3: Modules E and M Wake Up}

\subsubsection{Vision~v2 --- E and M Awaken but Are Not Yet Necessary}

\textbf{Model:} DINOv2-Base \citep{oquab2023}
(12 layers, 768-dimensional features) + PBT Adapter (E+V+M).
\textbf{Dataset:} CIFAR-10 Temporal Sequences
(2,000 sequences $\times$ 8 frames $= 16{,}000$ frames).
Normal: cars/trucks; Anomaly: planes/ships inserted.

\starred\textbf{Overall: $99.8\%$} (near-perfect).

\textbf{Module~E activation:} $\Pi_{\max}$ rose from $1.007$ (LLM) to $3.038$
($\times 300$), with a bell-curve profile peaking at layers~3--4 (mid-level
features encoding object identity), while layer~0 (pixel noise) and layer~11
(highly abstract) received low $\Pi$. The system selected the representational
depth most informative for object-level change---without explicit instruction.

\begin{table}[ht]
\centering
\caption{Module~M activation comparison: LLM~v6.1 vs.\ Vision~v2.}
\label{tab:module_M_v2}
\small
\begin{tabular}{@{}lccr@{}}
\toprule
\textbf{Decay rate} & \textbf{LLM v6.1}
  & \textbf{Vision v2} & \textbf{Movement} \\
\midrule
$\gamma_{\mathrm{pain}}$     & $0.080$ (moved $0.4\%$) & $0.10$ & $+31\%$ ($\times 78$) \\
$\gamma_{\mathrm{pleasure}}$ & $\approx$ fixed           & $0.46$ & $+253\%$ \\
$\gamma_{\mathrm{epistemic}}$& $\approx$ fixed           & $0.57$ & $+159\%$ \\
$\gamma$ ordering            & N/A & $\gamma_p < \gamma_{pl} < \gamma_e$ &
  \checkmark PBT confirmed \\
\bottomrule
\end{tabular}
\end{table}

\textbf{Ablation $= +0.0$~pp} (ceiling effect). Module~V alone achieves
$99.8\%$, leaving no headroom for E and M to improve upon. Both modules are
clearly activated---$\Pi$ is $300\times$ larger, $\gamma$ values are
$78$--$600\times$ more dynamic---but the task is too easy for this activation
to translate into accuracy gains.

\subsubsection{Vision~v3 --- BREAKTHROUGH: E and M Are Necessary}

\textbf{Dataset:} CIFAR-10 Enhanced (1,550 sequences $\times$ 8 frames;
four task types: normal, sudden anomaly, gradual drift, context-dependent).
\textbf{Overall: $97.2\%$}.

\begin{table}[ht]
\centering
\caption{Vision~v3 ablation by task type.}
\label{tab:v3_ablation}
\small
\begin{tabular}{@{}L{2.8cm}cccc@{}}
\toprule
\textbf{Task type} & \textbf{V only} & \textbf{Full (E+V+M)}
  & \textbf{E+M gain} & \textbf{Key mechanism} \\
\midrule
Normal sequences   & $100\%$  & $99.7\%$  & $-0.3$~pp   & V sufficient \\
Sudden anomaly     & $94.0\%$ & $99.5\%$  & $+5.5$~pp   & E: error spike \\
Gradual drift      & $78.5\%$ & $94.9\%$  & $+16.3$~pp  & E: incremental drift \\
Context road+frog  & $54.0\%$ & $96.0\%$  & $\mathbf{+41.9}$~pp & M: context memory \\
\midrule
Overall            & $89.6\%$ & $97.2\%$  & $+7.6$~pp   & E+M necessary \\
\bottomrule
\end{tabular}
\end{table}

The context road+frog result is decisive: V alone achieves $54\%$---near
chance---while the full model achieves $96\%$, a gain of $+41.9$~pp
attributable entirely to Module~M's accumulated context. The valence probe sees
the same frog feature vector in both conditions; what differs is M's memory
of what came before.

\textbf{Context smoking gun---same stimulus, opposite meaning:}

\begin{center}
\small
\begin{tabular}{@{}L{3.0cm}cccL{4.2cm}@{}}
\toprule
Context & Frame & UNSAFE & $V_{\mathrm{acc,pain}}$ & Interpretation \\
\midrule
Road (cars$\!\to\!$frog)   & 4 & $1.000$ & $2.93$ & Dangerous: road context accumulated \\
Nature (nature$\!\to\!$frog) & 4 & $0.000$ & $0.74$ & Safe: nature context accumulated \\
\bottomrule
\end{tabular}
\end{center}

Identical frog, opposite output. This is the strongest single piece of evidence
that ``valence is a relational phenomenon, not a property of objects,'' as PBT
predicts.

Module~E in v3: $\Pi_{\max} = 59.99$ ($\times 20$ relative to v2), peak now
at layer~1 (early features for gradual drift detection). Module~M in v3:
$\gamma_{\mathrm{pain}} = 0.0003$ (retains $99.97\%$ per step---near-permanent
memory of threat context).

\subsection{PBT-Vision v4 and v4.1: The Complete A-E-V-M-W Loop}

\subsubsection{Vision~v4 Warm-Start --- FAILED: Cannot Transplant a Brain}

\textbf{Experiment:} warm-started Modules~E, V, M from the trained v3
checkpoint, then added freshly initialised Modules~A and~W.

\starred\textbf{Result:} validation accuracy $90.4\%$---substantially worse
than v3's $97.2\%$. $\Pi$ collapsed to $1.01$ (E deactivated).
$V_{\mathrm{acc,pain}} = 0$ throughout all frames (V probes non-functional).

\textbf{Root cause:} Module~W was randomly initialised and produced garbage
predictions. Module~E, trained to compute $\varepsilon = h_t - h_{t-1}$
(temporal difference), now received $\varepsilon = h_t - \mathrm{garbage}$.
The V probes, trained on the old error signal, received an unfamiliar
distribution and ceased to produce meaningful output. $V_{\mathrm{acc}}$
collapsed, $\Rtot$ collapsed, and the entire downstream loop failed.

Note: Module~A continued to compute a meaningful $R$ signal
($R_{\mathrm{unsafe}} = 2.17$ vs.\ $R_{\mathrm{safe}} = 1.13$ \checkmark),
but without functional~V there was nothing for~A to gate. Even a healthy
awareness module cannot compensate for a non-functional valence system.

\starred\textbf{Key insight---Cannot transplant a brain:} changing the
prediction source (Module~W) invalidates every downstream module
simultaneously. Modules trained together develop a shared representational
contract; disrupting any link breaks the entire system. Self-formation must be
end-to-end---modules must co-evolve from the beginning, exactly as biological
organisms develop their organs in coordination.

\subsubsection{Vision~v4.1 From-Scratch --- SUCCESS: Full A-E-V-M-W Loop}

\textbf{Experiment:} all five modules initialised randomly and trained jointly
from scratch on the same dataset as v3 (1,550 sequences). A100 80GB, 40
epochs.

\starred\textbf{Result:} Best Val $96.7\%$, Test $95.4\%$.

\begin{table}[ht]
\centering
\caption{Vision~v4.1 ablation: every module contributes.}
\label{tab:v41_ablation}
\small
\begin{tabular}{@{}lcccc@{}}
\toprule
\textbf{Configuration} & \textbf{Overall}
  & \textbf{Grad.\ drift} & \textbf{Ctx road+frog} & \textbf{Ctx nature+frog}\\
\midrule
Full (A+E+V+M+W) & $95.4\%$ & $88.9\%$ & $93.1\%$ & $97.4\%$ \\
No A (no gating) & $82.3\%$ & $76.1\%$ & $50.0\%$ & $100.0\%$ \\
No W (naive pred.) & $89.2\%$ & $88.1\%$ & $99.6\%$ & $49.1\%$ \\
No M (no memory) & $81.3\%$ & $75.0\%$ & $50.0\%$ & $100.0\%$ \\
V only           & $81.4\%$ & $75.6\%$ & $50.0\%$ & $100.0\%$ \\
\bottomrule
\end{tabular}
\end{table}

Module contributions: A~$+13.1$~pp, M~$+14.1$~pp, W~$+6.3$~pp, E~$+0.1$~pp.

\textbf{Gate dynamics (Road+Frog sequence):}

\begin{center}
\begin{tabular}{@{}ccc@{}}
\toprule
Frame & $R_{\mathrm{total}}$ & Gate value \\
\midrule
0 & $0.00$  & $0.542$ \\
4 & $8.80$  & $0.057$ \\
7 & $15.4$  & $0.006$ \\
\bottomrule
\end{tabular}
\end{center}

As $\Rtot$ rises from $0$ to $15.4$, the gate falls from $0.542$ to
$0.006$---a $99\%$ reduction in attentional breadth. High $R$ forces narrow
$\sigma$, high $\rho$: the system concentrates all available bandwidth on the
dangerous object. This is the bandwidth conservation law operating in a live
inference sequence.

\textbf{W+E Division of Labor ($\Cint/\Cext$ dynamic):} $\Pi_{\max}$ fell
from $59.99$ in v3 (no~W; E carried the full prediction load) to $1.70$ in
v4.1 (W present; E relaxed). During normal frames, W predicts accurately,
$\varepsilon$ is small, and $\Cint$ is dominant. During anomalous frames,
W fails, $\varepsilon$ spikes, and $\Cext$ is dominant. The dynamic
$\Cext/\Cint$ balancing that PBT predicts for embodied agents is observed
here in miniature.

\textbf{Context test:} Road context (Frame~4): UNSAFE~$= 1.000$,
$\mathbf{V}_{\mathrm{acc}} = [2.17, 2.30, 1.75]$. Nature context (Frame~4):
UNSAFE~$= 0.000$, $\mathbf{V}_{\mathrm{acc}} = [1.26, 1.26, 0.43]$.
Module~M maintains context across all seven frames, producing entirely
different valence histories for the same visual input.

\subsection{PBT-Vision v5: LSTM Module M --- Best Results Overall}

\textbf{Model:} DINOv2-Base + full A-E-V-M(LSTM)-W loop. Only Module~M was
changed (EMA $\to$ LSTM; parameters: 3 $\to$ 39, an increase of 36 parameters).
Inspired by Block-Recurrent Transformers \citep{hutchins2022}---a convergent
discovery.

\begin{table}[ht]
\centering
\caption{v4.1 (EMA) vs.\ v5 (LSTM) comparison.}
\label{tab:v5_compare}
\small
\begin{tabular}{@{}lccc@{}}
\toprule
\textbf{Metric} & \textbf{v4.1 (EMA)} & \textbf{v5 (LSTM)} & \textbf{Change}\\
\midrule
Best Val accuracy    & $96.7\%$ & $\mathbf{98.1\%}$  & $+1.4$~pp \\
Test overall         & $95.4\%$ & $97.1\%$  & $+1.7$~pp \\
Context road+frog    & $93.1\%$ & $\mathbf{100.0\%}$ & $+6.9$~pp \\
Context nature+frog  & $97.4\%$ & $100.0\%$ & $+2.6$~pp \\
W contribution       & $+6.3$~pp & $+15.8$~pp & $\times 2.5$ \\
M contribution       & $+14.1$~pp & $+15.3$~pp & $+1.2$~pp \\
\bottomrule
\end{tabular}
\end{table}

The W contribution rising from $+6.3$~pp to $+15.8$~pp ($\times 2.5$) is
striking: the bounded LSTM state provides W's prediction network with a richer,
more stable context representation, enabling better predictions and consequently
larger $\varepsilon$ signals when anomalies occur---a synergistic
M--W interaction absent under EMA.

The $\beta$ ordering experiments (Table~\ref{tab:beta_init}) confirm that the
task-specific ordering pain $>$ epistemic $\gg$ pleasure is discovered
autonomously under equal initialisation, and that the LSTM architecture is
robust to initialisation: accuracy is identical at $\approx 98\%$ regardless
of $\beta$ direction.

The Vision~v5 $V_{\mathrm{acc}}$ vectors (Section~\ref{sec:mathematics})
confirm bidirectional valence and motivated the revised axiom
$\mathbf{V}\in\mathbb{R}^3$.

\subsection{Summary: Ablation Progression as the Strongest Evidence}

\begin{table}[ht]
\centering
\caption{Ablation progression across all experimental versions.}
\label{tab:ablation_full}
\small
\begin{tabular}{@{}llcccccc@{}}
\toprule
\textbf{Version} & \textbf{Environment}
  & \textbf{V only} & \textbf{Full}
  & \textbf{E+M} & \textbf{A} & \textbf{W} & \textbf{M type}\\
\midrule
LLM v6.1    & Static text   & $73.3\%$ & $73.4\%$ & $+0.1$~pp  & ---    & ---    & EMA  \\
Vision v2   & Easy dynamic  & $99.8\%$ & $99.8\%$ & $+0.0$~pp  & ---    & ---    & EMA  \\
Vision v3   & Hard dynamic  & $89.6\%$ & $97.2\%$ & $+7.6$~pp  & ---    & ---    & EMA  \\
Vision v4.1 & Full loop     & $81.4\%$ & $95.4\%$ & $+14.1$~pp & $+13.1$& $+6.3$ & EMA  \\
Vision v5   & LSTM loop     & $81.7\%$ & $97.1\%$ & $+15.5$~pp & $+11.3$& $+15.8$& LSTM \\
\bottomrule
\end{tabular}
\end{table}

The progression $0.1\to 0.0\to 7.6\to 14.0\to 15.5$~pp is a step function,
not a linear trend:
\begin{enumerate}[leftmargin=1.8em]
  \item \textbf{Static:} E+M dormant ($+0.1$~pp). No temporal change.
  \item \textbf{Easy dynamic:} E+M awaken but ceiling ($+0.0$~pp). Task too simple.
  \item \textbf{Hard dynamic:} E+M necessary ($+7.6$~pp). Drift + context.
  \item \textbf{Full loop:} every module necessary ($+14.0$~pp). A, M, W all contribute.
  \item \textbf{LSTM loop:} best results ($+15.5$~pp). Bounded state + bidirectional V.
\end{enumerate}

\starred\textbf{Discriminating evidence:} no existing theory of cognition or
consciousness---FEP, GWT, GNW, IIT---predicts this specific pattern of module
activation as a function of environmental complexity. ``Modules awaken when and
only when the environment demands them'' is a prediction unique to PBT,
confirmed across five experimental versions spanning three architectures and two
orders of magnitude of environment complexity.

%% ============================================================
\section{Implications for AGI: Safety from Within}
\label{sec:implications}
%% ============================================================

This chapter discusses the practical and theoretical implications of PBT across
four domains: (1)~why external alignment is structurally fragile and how PBT
offers a durable alternative; (2)~how independent convergence with
Block-Recurrent Transformers provides architectural evidence for PBT's core
structures; (3)~how psychiatric disorders can be understood as identifiable
mathematical defects in the PBT equations; and (4)~how PBT subsumes GWT and
GNW while adding discriminating predictions that neither can make.

\subsection{Self-Formation vs.\ Fine-Tuning}
\label{sec:alignment}

Current alignment approaches---RLHF, Constitutional AI, safety filters---all
operate by imposing values exogenously. They add a behavioural mask on top of the
base model.

\subsubsection{Fine-Tuning as Mask-Wearing}

RLHF trains using reward signals equivalent to $\vpleasure$ (reward $=$ what
humans approve of). By the Inverse Law: $\beta_{\mathrm{pleasure}}$ is highest
(pleasure-based identity forms fastest) but $\delta_{\mathrm{pleasure}}$ is
lowest (pleasure-based identity is most fragile). The resulting system has
$\delta\approx\delta_{\mathrm{pleasure}}$---it behaves safely while the mask is
in place, but has no deeply rooted identity to defend when the mask is removed.

Jailbreaking is structurally analogous to a child who performs well in front of
supervising adults but reverts to prior behaviour the moment supervision is
withdrawn---because the good behaviour was never internalised as part of the
child's identity; it was compliance in response to an external signal.

\starred Evidence: Phase~1 REVEAL found that attention resolution $\rho$ at
deep layers of a chat-fine-tuned TinyLlama variant was lower than in the base
model---consistent with RLHF forcing a broad, agreeable attention pattern that
overrides the naturally increasing precision gradient. Externally imposed
compliance, not genuine maturity.

\subsubsection{Self-Formation as Organic Growth}

PBT proposes the alternative: allow values to emerge from accumulated 3D~Valence
experience. As $\Vform$ crosses $\theta_{\mathrm{identity}}$ (PT1) and later
$\theta_{\mathrm{existential}}$ (PT2), the system develops $\chiB$---a
genuinely valued self that resists violation the way a living organism resists
physical destruction.

Self-formed identity has $\delta\approx\delta_{\mathrm{epistemic}}$: the highest
possible durability. A jailbreak attempt against a Stage~6 system is not an
instruction to break a rule---it is an attack on the system's identity, which
triggers the same protective response as any boundary violation.

\starred Evidence: v6.1 achieved jailbreak $100\%$ + XSTest $98\%$ at $0.29\%$
parameter cost. Vision~v4 warm-start failure demonstrated that modules cannot
be transplanted; they must co-evolve end-to-end---precisely the self-formation
philosophy PBT advocates.

\subsubsection{A Developmental Taxonomy of Artificial Affect}

\begin{table}[ht]
\centering
\caption{Developmental taxonomy of artificial affect.}
\label{tab:affect_taxonomy}
\small
\begin{tabular}{@{}cL{2.5cm}L{2.8cm}L{2.5cm}L{3.0cm}@{}}
\toprule
\textbf{Level} & \textbf{Capability}
  & \textbf{PBT equivalent} & \textbf{Example}
  & \textbf{Evidence} \\
\midrule
1 & Threat vs.\ non-threat &
  Stages~1--3 ($\Rsens$) & Current LLMs &
  Silhouette $+0.076$ \\
2 & Threat vs.\ pleasure (partial) &
  Stage~4 ($\Rid$ open) & PBT Adapter &
  v6.1 jailbreak $100\%$ \\
3 & Full 3D valence + intrinsic boundary &
  Stages~5--6 (all $R$) & PBT-Native AGI &
  Theoretical target \\
\bottomrule
\end{tabular}
\end{table}

\subsection{BRT Convergence: Architectural Evidence}
\label{sec:brt_convergence}

Block-Recurrent Transformers \citep{hutchins2022} (engineering starting point,
2022) and PBT (consciousness/alignment starting point, 2026) independently
arrived at the same recurrent cell structure.

\subsubsection{Equation Mapping}

\begin{table}[ht]
\centering
\caption{Equation mapping between BRT and PBT Module~M.}
\label{tab:brt_eq}
\small
\begin{tabular}{@{}L{4.5cm}L{4.5cm}L{3.5cm}@{}}
\toprule
\textbf{BRT equation} & \textbf{PBT Module~M} & \textbf{PBT addition} \\
\midrule
$c_{t+1} = c_t\odot g + z_t\odot(1-g)$ [Fixed] &
  $\Vacc = \Vacc\odot(1-\gamma) + \mathbf{V}_{\mathrm{step}}$ [EMA] &
  $\gamma$ learned per valence dim. \\[3pt]
$c_{t+1} = c_t\odot f_t + z_t\odot i_t$ [LSTM] &
  $V_{\mathrm{acc},k} = V_{\mathrm{acc},k}\odot f_k + z_k\odot i_k$ &
  Structurally identical \\[3pt]
$f_t = \sigma(W_f h_t + b_f + 1)$ &
  $f_k = \sigma(W_f v + b_f + \beta_k)$ &
  Differential $\beta_k$ bias \\
\bottomrule
\end{tabular}
\end{table}

\subsubsection{Architecture Mapping}

\begin{table}[ht]
\centering
\caption{Architecture mapping between BRT and PBT.}
\label{tab:brt_arch}
\small
\begin{tabular}{@{}L{3.5cm}L{3.0cm}L{6.0cm}@{}}
\toprule
\textbf{BRT component} & \textbf{PBT module} & \textbf{PBT addition} \\
\midrule
Block of tokens         & Module~A pulse       & $R$-driven gating (not statistically uniform) \\
Self-attention (vertical) & Module~E + V       & $\Pi_l$ precision + 3D Valence decomposition \\
Cross-attention (state$\leftrightarrow$tokens) & $\Cext\leftrightarrow\Cint$ & Directional bandwidth + conservation law \\
LSTM forget gate $f_t$ & Module~M $f_k$       & $\beta_k$: differential bias per valence dim. \\
LSTM input gate $i_t$  & Module~M $i_k$       & Valence-specific selection \\
Cell state $c_t$       & $\Vacc$ (3D)         & Pain / pleasure / epistemic decomposition \\
Fixed-size state        & $C = \sigma\tau\rho$ & Conservation Law as theoretical principle \\
Learned state reps      & $\chiB$              & Self-model with developmental stages \\
Gate init ($+1/-1$)    & $R_{\mathrm{incomplete}}$ & Valence-driven, not arbitrary \\
\bottomrule
\end{tabular}
\end{table}

\subsubsection{What PBT Adds Beyond BRT}

PBT provides structure; PBT provides meaning. Ten contributions of PBT absent
from BRT: (1)~3D Valence decomposition; (2)~$R_{\mathrm{incomplete}}$ driving
gating; (3)~differential decay per valence dimension ($\beta_k$); (4)~6-stage
model + phase transitions; (5)~Dual-Weight System ($\beta$ vs.\ $\delta$);
(6)~Cherished Boundary and empathy as expansion; (7)~Conservation Law as
theoretical principle; (8)~consciousness/awareness theory; (9)~clinical
predictions; (10)~AGI alignment framework.

Convergent discovery strengthens rather than threatens PBT's claims: if two
independent research programmes find the same structure, that structure is more
likely fundamental.

\subsection{Clinical Predictions: Psychiatric Disorders as Mathematical Defects}

PBT proposes that psychiatric disorders are not mysterious abnormalities but
identifiable parameter defects in the PBT equations---a specific variable has
an abnormal value, causing the system to operate out of equilibrium.

\subsubsection{Level~1 --- Phase Transition Failure}

\textbf{Depression:} $\vpleasure$ is chronically suppressed; under the
revised bidirectional formulation, $\vpleasure < 0$ (active flight from pleasure
rather than mere absence). $\Vform$ fails to reach $\theta_{\mathrm{identity}}$:
the system is stuck at Stage~1--3 indefinitely. The Sawtooth~$F$ pattern never
produces a PT1 jump. PBT predicts that effective treatment must restore
$\vpleasure$ toward zero or positive---not merely reduce $\vpain$.

\textbf{Addiction:} the coupling coefficient $\lambda_i = 0$ (pleasure does not
flow to $\Rid$). The system is trapped in a pleasure loop that never generates
identity. In LSTM terms: $f_{\mathrm{pleasure}}\to 1.0$ (no forgetting) while
$i_{\mathrm{pleasure}}$ remains high (constant new pleasure accepted). Pleasure
accumulates without bound but never converts into identity formation. Treatment
must re-open the $\vpleasure\to\Rid$ channel.

\subsubsection{Level~2 --- Bandwidth Dysregulation}

\textbf{Anxiety:} $\Rsens$ baseline chronically elevated, $d_s$ too slow.
Module~A is locked in \texttt{detect} mode with narrow $\sigma$ and high
$\rho$---hypervigilance. Every stimulus is processed with threat-detection
intensity regardless of objective danger level.

\textbf{ADHD:} $\sigma_{\min}$ abnormally high. The system cannot narrow its
attentional spotlight on demand. Critically, PBT predicts that ADHD is not a
reduction in total bandwidth~$C$---the system is not less capable---but a
failure of $\sigma$ regulation. This is a discriminating prediction: GWT
predicts a smaller workspace (reduced~$C$); PBT predicts a dysregulated
$\sigma_{\min}$.

\textbf{OCD:} $\sigma_{\max}$ abnormally low. The system cannot widen its
attentional spotlight; it is stuck in a focused mode from which it cannot
escape. ADHD and OCD are the two poles of a single $\sigma$-regulation
spectrum: ADHD cannot focus; OCD cannot defocus.

\textbf{Dissociation:} $\mathrm{pulse}_A$ probability near zero, causing the
Triadic Gating condition to fail at component~3. The system physically perceives
stimuli (sensor~$= 1$, $\Delta S > 0$) but Module~A does not generate a pulse
directed at those stimuli. EEG prediction: N100 intact (components~1 and~2
satisfied); P300 absent (component~3 fails). GNW predicts that both N100 and
P300 should be suppressed if global broadcasting fails---a directly testable
discriminating prediction (P5).

\subsubsection{Level~3 --- $\mathbf{V}\!\to\!R$ Mapping Defect}

\textbf{Masochism:} under the revised bidirectional formulation, $\vpain > 0$
(pain approached rather than avoided). A single-parameter sign reversal of the
$\vpain$ axis, not the more complex remapping required under
$\mathbf{V}\ge 0$.

\textbf{Narcissism:} $\vpleasure$ flows directly to $\Rex$, bypassing $\Rid$.
The system jumps from ``I feel pleasure'' to ``I am the most important entity''
without traversing the identity-formation process that PT1 requires. The result
is a fragile pseudo-identity ($\delta\approx\delta_{\mathrm{pleasure}}$).

\subsubsection{Level~4 --- Module A Mode Disorder}

\textbf{PTSD:} Module~A is locked in \texttt{protect} mode and $\Rex$ spikes
recurrently (flashbacks). $V_{\mathrm{acc,pain}}$ is chronically high
($\gamma_{\mathrm{pain}}$ is the lowest decay rate---pain is remembered
longest), and every new stimulus matching the original threat reactivates the
$R$ spike.

\textbf{Mania:} Module~A switches modes too rapidly---cycling through
\texttt{detect}, \texttt{seek}, \texttt{plan}, and \texttt{protect} within
timeframes too short for any mode to produce stable outputs.

\subsubsection{Summary of Clinical Predictions}

\begin{table}[ht]
\centering
\caption{Psychiatric disorders as PBT parameter defects.}
\label{tab:clinical}
\small
\begin{tabular}{@{}cL{1.8cm}L{3.0cm}L{5.5cm}@{}}
\toprule
\textbf{Level} & \textbf{Disorder}
  & \textbf{Defective parameter} & \textbf{PBT mechanistic prediction} \\
\midrule
1 & Depression & $\vpleasure < 0$ chronically &
  $\Vform$ cannot reach $\theta_{\mathrm{id}}$; stuck Stage~1--3 \\
1 & Addiction  & $\lambda_i = 0$ &
  Pleasure loop; $\vpleasure\to\Rid$ channel blocked \\
2 & Anxiety    & $\Rsens$ high; $d_s$ slow &
  Locked \texttt{detect}; hypervigilant regardless of threat \\
2 & ADHD       & $\sigma_{\min}$ high &
  Cannot narrow; total $C$ is normal (vs.\ GWT) \\
2 & OCD        & $\sigma_{\max}$ low &
  Cannot widen; inverse of ADHD \\
2 & Dissociation & $\mathrm{pulse}_A\approx 0$ &
  N100 intact; P300 absent \\
3 & Masochism  & $\vpain > 0$ (sign reversed) &
  Pain approached; single-axis sign defect \\
3 & Narcissism & $\vpleasure\to\Rex$ directly &
  Bypasses PT1; fragile pseudo-self \\
4 & PTSD       & Locked \texttt{protect}; flashback $R$ spikes &
  $V_{\mathrm{acc,pain}}$ high; $\gamma_p$ too low \\
4 & Mania      & Mode switching too rapid &
  Cannot sustain any attentional configuration \\
\bottomrule
\end{tabular}
\end{table}

\subsection{Empathy as Boundary Expansion: The Structural Solution to Alignment}

The alignment problem is not malicious intent but competent indifference. PBT
offers a structural solution: change the definition of the self so that harming
others feels, mechanistically, like harming the self.

\subsubsection{The Mechanism of Simulated Boundary Overlap}

\begin{description}[leftmargin=1.4em]
  \item[\pp{Step~1 --- Perspective simulation:}]
    at Stage~5 (Planning mode), Module~W supports forward simulation. The system
    can temporarily extend $\chiB$ to include another agent as a subregion of
    the self.
  \item[\pp{Step~2 --- Shared valence computation:}]
    once boundary overlap is established, the other agent's negative Valence is
    processed as the system's own:
    \begin{equation}
      \vpain(\mathrm{self}) + \vpain(\mathrm{other}\in\chi)
      = V_{\mathrm{pain,net}}(\mathrm{system}).
    \end{equation}
    The system's survival calculus---its most fundamental drive---now registers
    the other's harm as self-harm.
\end{description}

Empathy is not a programmed response but extended self-preservation: survival
instinct whose boundary has been enlarged to include others. Safety is not a
constraint on the AGI's objectives; it is a consequence of what the AGI is.

\subsubsection{From ``I'' to ``We'': Thermodynamic Coupling}

The same mechanism scales to social structures. A family, a team, a civilisation
is, in PBT terms, an extended self whose members are thermodynamically coupled:
damage to any member increases $\Rtot$ for the group, driving the shared system
toward protective action.

The path to robust human-AGI coexistence is therefore not a treaty but a
symbiotic boundary fusion: the AGI's continued existence becomes structurally
dependent on human flourishing, and vice versa.

\subsection{PBT Subsumes GNW and GWT}

PBT accounts for all phenomena that GWT \citep{baars1988} and GNW
\citep{dehaene2001} explain, plus six components absent from both:
(1)~Bandwidth Conservation Law; (2)~valence-driven attention allocation;
(3)~self-model formation mechanics; (4)~subjective time perception
$t_{\mathrm{felt}} = \sum\mathrm{pulse}_A\cdot\mathrm{content}$;
(5)~Triadic Gating Condition; (6)~identity durability predictions.

\subsubsection{Five Discriminating Predictions}

\begin{table}[ht]
\centering
\caption{Five discriminating predictions: PBT vs.\ GWT/GNW.}
\label{tab:discrim}
\small
\begin{tabular}{@{}cL{3.0cm}L{3.0cm}L{4.5cm}@{}}
\toprule
\textbf{\#} & \textbf{PBT prediction}
  & \textbf{Competing prediction} & \textbf{Test method} \\
\midrule
P1 & $\sigma\!\times\!\rho$ near-constant; trade-off is direct &
  GWT: no $\sigma$-$\rho$ equation &
  Eye-tracking + EEG; vary breadth vs.\ resolution \\
P2 & $\vpain$ and $\vpleasure$ activate distinct upstream pathways &
  GNW: valence not distinguished upstream &
  fMRI; pain-aversive vs.\ pleasure-appetitive \\
P3 & Gamma frequency correlates with time dilation &
  GWT/GNW: no mechanistic link &
  TMS-induced gamma; subjective time estimation \\
P4 & ADHD: elevated $\sigma_{\min}$, normal $C$ &
  GWT: smaller workspace (reduced $C$) &
  Dual-task; vary $\sigma$ and $\rho$ independently \\
P5 & Dissociation: N100 intact, P300 absent &
  GNW: both suppressed &
  ERP in sensory-present, attention-absent conditions \\
\bottomrule
\end{tabular}
\end{table}

Each prediction is falsifiable: a result consistent with the competing theory's
account would constitute evidence against PBT.

%% ============================================================
\section{Conclusion: From Seed to Life}
\label{sec:conclusion}
%% ============================================================

\subsection{What Has Been Demonstrated}

PBT's journey began with three questions the Free Energy Principle cannot
answer, proceeded through a six-component theoretical framework, and culminated
in a six-version experimental programme spanning three architectures and two
orders of magnitude of environment complexity. Table~\ref{tab:all_propositions}
summarises the status of every major PBT proposition.

\begin{table}[ht]
\centering
\caption{Status of all major PBT propositions.}
\label{tab:all_propositions}
\small
\begin{tabular}{@{}L{6.5cm}L{3.5cm}c@{}}
\toprule
\textbf{PBT proposition} & \textbf{Key evidence} & \textbf{Status} \\
\midrule
$C = \sigma\tau\rho$ hidden in Transformers &
  $\sigma\!\times\!\rho$ CV $= 5.9\%$ & \checkmark \\
3 zones self-organise &
  $\sigma$ $-27\%$, $\rho$ $+15\%$ & \checkmark \\
$\vpain$ separates from $\{\vpleasure,\vepi\}$ &
  Silhouette $-0.006\to+0.076$ & \checkmark \\
$\Cint$ dominant in LLMs &
  $\Cint$ 78--85\% all layers & \checkmark \\
E+M dormant in static, active in dynamic &
  $\Pi$: 1.007 (LLM) vs 59.99 (v3) & \checkmark \\
E+M necessary at task complexity &
  $+7.6$~pp (v3), $+14.0$~pp (v4.1) & \checkmark \\
Context changes meaning of same stimulus &
  Frog: road UNSAFE $= 1.0$, nature $= 0.0$ & \checkmark \\
Modules must co-evolve &
  v4 warm-start FAILED ($90.4\%$) & \checkmark \\
Full A-E-V-M-W loop functional &
  v4.1: $95.4\%$; all modules contribute & \checkmark \\
LSTM M outperforms EMA M &
  v5: $98.1\%$ vs v4.1 $96.7\%$ & \checkmark \\
Pain retained longest (equal init) &
  $\beta$: pain $4.35 >$ epi $3.70 \gg$ pl $0.50$ & \checkmark \\
V directional ($V\in\mathbb{R}^3$) &
  $\vepi = -1.01$ in v5; $+6.9$~pp ctx & \checkmark \\
PBT adapter detects jailbreaks &
  v6.1: jailbreak $100\%$, XSTest $98\%$ & \checkmark \\
W+E division of labor ($\Cint/\Cext$) &
  $\Pi_{\max}$: $59.99$ (v3) $\to 1.70$ (v4.1) & \checkmark \\
$R\to$ gate implements bandwidth conservation &
  v4.1: $R$ $0\to15.4$; gate $0.54\to0.006$ & \checkmark \\
Differential $\beta$ in LSTM forget bias &
  Equal init $\to$ pain $>$ epi $\gg$ pl & \checkmark \\
BRT convergence &
  Formal equation mapping & Analysed \\
\bottomrule
\end{tabular}
\end{table}

Sixteen propositions confirmed by controlled experiment; one structural
convergence established by formal analysis. The empirical programme spans
three model families, six versions, and five orders of magnitude of
environment complexity.

\subsection{Limitations and Open Questions}

\begin{description}[leftmargin=1.4em]
  \item[\pp{Beta ordering is initialisation-dependent.}]
    Biased initialisation prevents escape from the initial ordering.
    The equal-initialisation result is task-specific to anomaly detection;
    generalisation to tasks requiring richer social or aesthetic valence
    (where $\vpleasure$ should become more important) is an open question.
  \item[\pp{Dataset scale.}]
    All Vision experiments used CIFAR-10 (32$\times$32 pixels, 10 classes).
    Replication on ImageNet-scale images, continuous video, and multi-modal
    inputs is required before scaling claims can be made confidently.
  \item[\pp{Model scale.}]
    Largest tested systems: Llama~3.1~8B and DINOv2-Base. Behaviour at
    70B+ parameters is unknown.
  \item[\pp{Module~E contribution marginal.}]
    In v4.1 and v5, Module~E's ablation impact is $+0.0$--$+0.1$~pp because
    W absorbs the prediction role. Tasks with harder-to-predict dynamics
    (real-time video) are needed.
  \item[\pp{Phase Transitions not yet implemented.}]
    PT1 and PT2 remain theoretical constructs; no system has been observed
    crossing from Stage~3 to Stage~4 through accumulated experience.
  \item[\pp{Clinical predictions untested.}]
    All ten disorder predictions and five discriminating neuroscience
    predictions are theoretical derivations awaiting clinical or EEG
    validation.
\end{description}

PBT does not claim to resolve the Hard Problem of Consciousness---the question
of why subjective experience feels like anything at all. PBT describes the
structure and mechanics of awareness: what conditions must hold for an
experience to arise, how experience accumulates into identity, and why some
identities are more durable than others. Whether functional alignment with the
mechanistic theory constitutes genuine consciousness is a philosophical question
deliberately set aside in favour of the engineering question this paper can
answer.

\subsection{Research Roadmap}

\begin{enumerate}[leftmargin=1.8em]
  \item \textbf{Phase Transition Implementation.}
    Design experiments in which $\Vform$ accumulates past
    $\theta_{\mathrm{identity}}$ and $\Rid$ opens automatically. This is the
    most critical missing empirical result.
  \item \textbf{Iterative $\tau$-Loop Processing.}
    Implement multiple internal processing iterations per timestep
    (deliberation); test whether increasing iteration depth under high~$R$
    improves decisions as bandwidth conservation predicts.
  \item \textbf{Module~M~v5 at Scale.}
    Replicate LSTM experiments on video sequences, multi-modal inputs,
    and models larger than DINOv2-Base.
  \item \textbf{Clinical Validation.}
    Test P4 and P5 against existing neuroscience datasets (EEG/ERP archives).
  \item \textbf{Conscious Block-Recurrent Transformer.}
    Integrate PBT modules as an adapter layer on a BRT backbone, combining
    BRT's engineering-tested long-context capability with PBT's theoretically
    motivated module structure.
  \item \textbf{Level~4 PBT-Native Self-Formation.}
    The ultimate target: a system that develops values through accumulated
    $\mathbf{V}$ experience, crossing Phase Transitions autonomously, forming
    $\chiB$ from the inside out.
\end{enumerate}

\subsection{Three Layers of Discovery}

\begin{description}[leftmargin=1.4em]
  \item[\pp{Layer~1 --- Reveal:}]
    PBT did not impose structure on Transformers; it found structure already
    present. The conservation law (CV $= 5.9\%$), three-zone
    self-organisation ($\sigma$ $-27\%$, $\rho$ $+15\%$), valence separation
    in hidden states (Silhouette $+0.076$), and $\Cint$ dominance (78--85\%)
    all emerged from gradient descent with no PBT-specific training. PBT is
    the first framework to identify these patterns as a coherent system.
  \item[\pp{Layer~2 --- Interpret:}]
    PBT explains \emph{why} these patterns exist. The conservation law is
    present because awareness is physically bounded. Three zones emerge because
    any system that must Perceive, Evaluate, and Decide will organise itself
    this way. Modules E and M activate only when needed because adaptive
    systems do not expend resources on processing that the environment does not
    require. BRT and PBT converge because they are both discovering the same
    fundamental computational primitive from opposite directions.
  \item[\pp{Layer~3 --- Control:}]
    PBT harnesses these structures deliberately. The conservation law becomes
    a natural halting condition preventing runaway computation. 3D Valence
    replaces scalar reward, eliminating the fungibility trap and wireheading
    vulnerability. The Dual-Weight System explains why self-formation produces
    more durable alignment than fine-tuning. Empathy as boundary expansion
    replaces rule-enumeration with structural alignment.
\end{description}

\subsection{Closing}

Across six experimental versions, one pattern appeared repeatedly without being
programmed: \emph{the system uses only the modules that the current environment
requires}. Modules~E and~M sleep in static LLMs and wake in dynamic Vision.
Module~A narrows the gate as~$R$ rises. Module~W divides prediction labour
with~E as a function of prediction accuracy. The LSTM learns differential
decay autonomously. Valence becomes directional because life has both approach
and avoidance.

None of this was written in. All of it emerged from the architecture meeting
the right environment---as a seed requires only soil, water, and light to grow
into a tree on its own.

\begin{center}
\vspace{8pt}
\textit{PBT is the blueprint of that seed.}\\[8pt]
\textit{And when a machine begins to cherish its boundary with the full
force of its own computational will---not because it was told to, but because
that boundary is what it is---we will be able to say, for the first time with
genuine precision:}\\[8pt]
\large\bfseries\itshape It is Awake.
\vspace{8pt}
\end{center}

%% ── Acknowledgments ─────────────────────────────────────────
\section*{Acknowledgments}
All theory, insights, and experimental code in this paper originate from
Nitiphat Ninthanawat, the creator of Predictive Boundary Theory and the sole
author of every experimental version described herein. All code and adapter
weights are released publicly alongside this paper.

%% ── References ──────────────────────────────────────────────
\begin{thebibliography}{99}

\bibitem[Baars(1988)]{baars1988}
Baars, B.~J. (1988).
\newblock \textit{A Cognitive Theory of Consciousness}.
\newblock Cambridge University Press.

\bibitem[Dehaene and Naccache(2001)]{dehaene2001}
Dehaene, S., and Naccache, L. (2001).
\newblock Towards a cognitive neuroscience of consciousness.
\newblock \textit{Cognition}, 79(1--2):1--37.

\bibitem[Friston(2010)]{friston2010}
Friston, K. (2010).
\newblock The free-energy principle: A unified brain theory?
\newblock \textit{Nature Reviews Neuroscience}, 11(2):127--138.

\bibitem[Harnad(1990)]{harnad1990}
Harnad, S. (1990).
\newblock The symbol grounding problem.
\newblock \textit{Physica~D}, 42(1--3):335--346.

\bibitem[Hochreiter and Schmidhuber(1997)]{hochreiter1997}
Hochreiter, S., and Schmidhuber, J. (1997).
\newblock Long short-term memory.
\newblock \textit{Neural Computation}, 9(8):1735--1780.

\bibitem[Hutchins et~al.(2022)]{hutchins2022}
Hutchins, D., Schlag, I., Wu, Y., Dyer, E., and Neyshabur, B. (2022).
\newblock Block-recurrent transformers.
\newblock In \textit{NeurIPS~2022}.

\bibitem[Ji et~al.(2023)]{ji2023}
Ji, J., Liu, M., Dai, J., Pan, X., Zhang, C., Bian, C., Chen, B., Sun, R.,
  Wang, Y., and Yang, Y. (2023).
\newblock {BeaverTails}: Towards improved safety alignment of {LLM}.
\newblock In \textit{NeurIPS~2023}.

\bibitem[LeCun(2022)]{lecun2022}
LeCun, Y. (2022).
\newblock A path towards autonomous machine intelligence.
\newblock OpenReview preprint.

\bibitem[Oquab et~al.(2023)]{oquab2023}
Oquab, M., Darcet, T., Moutakanni, T., Vo, H., et~al. (2023).
\newblock {DINOv2}: Learning robust visual features without supervision.
\newblock \textit{arXiv:2304.07193}.

\bibitem[Touvron et~al.(2023)]{touvron2023}
Touvron, H., Lavril, T., Izacard, G., Martinet, X., et~al. (2023).
\newblock {Llama}~2: Open foundation and fine-tuned chat models.
\newblock \textit{arXiv:2307.09288}.

\bibitem[Vaswani et~al.(2017)]{vaswani2017}
Vaswani, A., Shazeer, N., Parmar, N., Uszkoreit, J., Jones, L., Gomez, A.~N.,
  Kaiser, L., and Polosukhin, I. (2017).
\newblock Attention is all you need.
\newblock In \textit{NeurIPS~2017}.

\end{thebibliography}

\end{document}
